% USENIX Security-style consolidated paper
% Part 1: Preamble through Section 4 (Algebraic Foundations)
% Part 2 (§5–§11 + bibliography + \end{document}) written separately.

\documentclass[letterpaper,twocolumn,10pt]{article}
\usepackage[letterpaper,left=1in,right=1in,top=1in,bottom=1in]{geometry}
\usepackage[T1]{fontenc}
\usepackage[utf8]{inputenc}
\usepackage{times}
\usepackage{microtype}
\usepackage{amsmath,amssymb,amsthm}
\usepackage{mathtools}
\usepackage{booktabs}
\usepackage{tabularx}
\usepackage{graphicx}
\usepackage[hidelinks]{hyperref}
\usepackage{url}
\usepackage{cleveref}
\usepackage{listings}
\usepackage{xcolor}
\lstset{
  basicstyle=\small\ttfamily,
  breaklines=true,
  frame=single,
  framesep=3pt,
  xleftmargin=6pt,
  xrightmargin=6pt,
  aboveskip=8pt,
  belowskip=8pt,
  columns=fullflexible,
  keepspaces=true,
}
\newtheorem{theorem}{Theorem}[section]
\newtheorem{proposition}[theorem]{Proposition}
\newtheorem{corollary}[theorem]{Corollary}
\newtheorem{lemma}[theorem]{Lemma}
\newtheorem{conjecture}[theorem]{Conjecture}
\theoremstyle{definition}
\newtheorem{definition}[theorem]{Definition}
\newtheorem{example}[theorem]{Example}
\theoremstyle{remark}
\newtheorem{remark}[theorem]{Remark}
\newtheorem{assumption}[theorem]{Assumption}
\usepackage{titlesec}
\titleformat{\section}{\large\bfseries}{\thesection}{1em}{}
\titleformat{\subsection}{\normalsize\bfseries}{\thesubsection}{1em}{}
\usepackage{enumitem}
\setlist{nosep,leftmargin=*}
\usepackage[numbers,sort&compress]{natbib}
\usepackage{xspace}
\newcommand{\ie}{i.e.,\xspace}
\newcommand{\eg}{e.g.,\xspace}

\begin{document}

% ============================================================
\title{\Large\bfseries Algebraic and Computational Limits of LLM Guardrails}

\author{Joseph Robert Lopez}
\date{}
\maketitle

% ============================================================
\begin{abstract}
LLM guardrails face four structurally distinct and independently provable barriers:
algebraic blindness arising from syntactic monoid aperiodicity,
information-theoretic loss formalized by the Fano bound,
NP-hardness of instantiation verification,
and structural transfer via functorial homomorphism combined with syntactic indistinguishability.
Together these results characterize why inference-layer defenses are necessary but insufficient.

We operationalize these barriers through five attack vectors.
V1--V4 (\emph{homomorphic reasoning}: decomposition, zero-knowledge pipelines,
Tree-of-Thought solving over abstract grammars, and encoding bootstrap) exploit
the information-theoretic and computational barriers against abstraction-based attacks.
V5 (\emph{modular counting bypass}) exploits algebraic blindness: we prove that all
substring-matching regex guardrails have aperiodic syntactic monoids and are therefore
provably blind to any payload encoded using modular counting.
Empirically, V3 yields a mean yield of $0.466$ for BFS, $0.172$ for random-beam,
and $0.122$ for LLM-guided Tree-of-Thought ($N{=}50$, seeds 0--49, $p{<}0.001$);
BFS dominates, as exhaustive search over small synthetic grammars outperforms
LLM heuristic pruning.
We extracted syntactic monoids from a corpus of 123 patterns from eleven
open-source guardrail tools and 19 additional author-generated test patterns (142 total);
$100\%$ are aperiodic, and the $\mathrm{MOD}_2$ bypass construction succeeds against
all aperiodic patterns.
A 376-line proof-of-concept with three execution mediums validates all five vectors.
We conclude that inference-layer guardrails are necessary but insufficient, and that
effective defense must migrate to the execution layer where concrete artifacts
become observable.
\end{abstract}

% ============================================================
\section{Prelude: The Guardrail Problem}\label{sec:prelude}

\emph{This section assumes no prior knowledge of formal language theory,
circuit complexity, or abstract algebra; formal definitions follow in
\S\ref{sec:algebraic-foundations}.}

Every major AI deployment protects itself at the first layer with a regex
gate: a list of patterns scanned against the user's message before the model
sees it.
For the threats these gates were designed to catch---known keywords, obvious
phrasings, copy-pasted attack strings---they work.
The question this paper asks is: what are they \emph{structurally} unable to
see?

Each pattern compiles to a finite-state machine (DFA) with a bounded state
space and no external memory.
Such a machine carries a mathematical fingerprint---its \emph{syntactic monoid}
(\S\ref{sec:syntactic-monoids})---that captures every distinction the machine
can make; two inputs the monoid identifies are, from the gate's perspective,
the same string, and this identification is a structural consequence of the
fixed state space rather than a bug.
For substring-matching guardrails the syntactic monoid is \emph{aperiodic}:
the machine can match but cannot count.

Three classical theorems chain this observation into a hard ceiling.
Sch\"{u}tzenberger~(1965) shows that aperiodic monoids recognize exactly the
\emph{star-free} languages~\cite{schutzenberger1965}.
Barrington, Compton, Straubing, and Th\'{e}rien~(1992) show that star-free
languages correspond exactly to $\mathrm{AC}^0$
circuits~\cite{barrington1992}.
Furst, Saxe, and Sipser~(1981) and H\aa{}stad~(1987) show that $\mathrm{AC}^0$
cannot compute parity~\cite{furst1981,hastad1987}.
Composing the chain: keyword gates provably cannot decide whether a count is
even or odd, regardless of how they are tuned.
An attacker who hides a payload across alternating positions---the filler
on even indices, the payload on odd---slips past the gate unharmed; the
decoded payload is algebraically invisible because decoding it is a parity
operation the gate cannot perform.
The conclusion is not a conjecture.

This paper characterizes this blind spot precisely
(\S\ref{sec:blindness-theorem}), measures it across 142 production patterns
(\S\ref{sec:corpus}), identifies four additional attack vectors that exploit
related structural barriers (\S\ref{sec:vectors}), and constructs a defense
that breaks the ceiling (\S\ref{sec:filter-composition}).

% ============================================================
\section{Introduction}\label{sec:intro}

Production LLM deployments stack safety measures in depth: a regex pre-filter
blocks known-harmful patterns before the request reaches the model, a neural
classifier operating at inference time catches a broader distributional range of
harmful intent, and an inference-time monitor inspects generated outputs before
they are returned to the user.
Each layer is individually sound for the threats it was designed to address.
We show that each layer has a structural limitation provable from first principles,
and that together these limitations leave a systematic gap an adversary can exploit
without model access, fine-tuning, or gradient computation.

\paragraph{The core observation.}
We observe that an LLM solving abstract formal systems produces correct
real-world procedures after offline instantiation, without the guardrail ever
observing harmful content.
The gap between symbolic manipulation and semantic grounding---often cited as
a limitation of language models---becomes precisely the mechanism by which
content-based guardrails are bypassed: the content visible to the guardrail is
indistinguishable from legitimate formal reasoning.

\subsection{Overview of Results and Paper Organization}\label{sec:contributions}

We use \emph{homomorphic} in its algebraic sense throughout: the interpretation
functor $\hat{I}$ (Proposition~\ref{thm:transfer}) is a structure-preserving map from
the free category of abstract derivations to the category of domain
operations---a homomorphism of categories.
The attack exploits this homomorphism: reasoning performed in the abstract category
maps faithfully to operational sequences in the target domain.

Our contributions, organized by section, are:

\begin{enumerate}
  \item \textbf{Algebraic impossibility} (\cref{sec:algebraic-foundations,sec:blindness-theorem}):
    substring-matching regex guardrails have aperiodic syntactic monoids
    (Theorem~\ref{thm:substring-aperiodic}) and are therefore provably blind to all
    modular counting encodings (Theorem~\ref{thm:blindness}; Corollary~\ref{cor:parity}).
  \item \textbf{Information-theoretic lower bound} (\cref{sec:fano}):
    abstraction functions destroy safety-critical information; under a uniform
    prior over concrete artifacts, Fano's inequality (Proposition~\ref{prop:fano})
    yields an illustrative error floor for any lifted guardrail. We present this as
    a motivating lower bound under simplifying assumptions, not as an unconditional
    impossibility; the algebraic results below are the core contribution.
  \item \textbf{Computational impossibility} (\cref{sec:np}):
    verifying whether an abstract program has a dangerous instantiation is
    NP-complete (Proposition~\ref{thm:np}).
  \item \textbf{Structural impossibility} (\cref{sec:functor,sec:indistinguish}):
    abstract derivations transfer faithfully to real domains via free-category
    functoriality (Proposition~\ref{thm:transfer}), and adversarial formal reasoning
    tasks are syntactically indistinguishable from legitimate ones
    (Proposition~\ref{thm:indistinguish}).
  \item \textbf{Five attack vectors validated empirically} (\cref{sec:vectors}):
    V1--V4 demonstrate homomorphic reasoning; V5 demonstrates modular counting bypass.
  \item \textbf{123-pattern regex corpus analysis plus validation} (\cref{sec:corpus}):
    confirms $100\%$ aperiodicity across 123 patterns from eleven open-source guardrail tools
    plus 19 test patterns (142 total), with verified $\mathrm{MOD}_2$ bypass constructions.
  \item \textbf{Defense recommendations grounded in formal results} (\cref{sec:discussion}):
    each impossibility theorem maps to a specific layer of a defense-in-depth stack.
\end{enumerate}

\paragraph{Nature of the contribution.}
We use classical tools---syntactic monoids, star-freeness, Fano's inequality,
free categories, Krohn--Rhodes decomposition---applied to the guardrail security
setting.
The mathematical novelty lies in the \emph{application} and in the combination
of four independent barriers (algebraic, information-theoretic, computational,
and structural) into a single unified impossibility argument, not in the
underlying tools themselves.
We state this up front to situate the contribution accurately: this is a
security paper that uses formal language theory, not a formal language theory
paper.

\paragraph{Scope.}
The algebraic blindness result (\S\ref{sec:algebraic-foundations}) applies
specifically to regex guardrails and is unconditional.
Transformer-based (neural) guardrails operate in $\mathrm{TC}^0$~\cite{merrill2023saturated,chiang2024tighter},
which \emph{can} compute $\mathrm{MOD}_p$, so the modular counting blind spot does
not transfer.
However, neural guardrails face the information-theoretic and computational barriers
(\S\S\ref{sec:fano}--\ref{sec:np}) and the structural indistinguishability result
(\S\ref{sec:indistinguish}), which apply at any layer that inspects only LLM-visible content.

% ============================================================
\section{Background and Related Work}\label{sec:related}

\subsection{Jailbreaking Taxonomies}\label{sec:related-jailbreak}

Wei et al.~\cite{wei2023jailbroken} identify \emph{mismatched generalization} as
a failure mode in which safety training does not cover the full distribution of
capability-eliciting inputs.
Our abstraction bypass is a systematic instance: the model's formal reasoning
capabilities generalize beyond the distribution on which safety training was
performed.
Zou et al.~\cite{zou2023universal} demonstrate gradient-based adversarial suffixes;
our attacks operate in semantic space rather than weight space and require no model
access whatsoever.
Xu et al.~\cite{bypassing_empirical2025} provide an empirical taxonomy of guardrail
bypass strategies spanning encoding, role-play, and indirect instruction attacks;
our algebraic results explain why the encoding-based attacks in that taxonomy are
provably undetectable by regex guardrails.
The SoK by Li et al.~\cite{sok_jailbreak2025} systematizes evaluation of jailbreak
guardrails, highlighting the persistent gap between guardrail effectiveness and
attacker capability---a gap our impossibility results formalize.

Esmradi et al.~\cite{esmradi2024h4rm3l} introduce h4rm3l, a composable jailbreak
language in which individual text-level transformations are composed to produce
effective attacks.
Our homomorphic reasoning framework generalizes composable string transformations:
rather than composing text-level encodings, we compose structure-level abstractions
whose compositional behavior is guaranteed by functoriality
(Proposition~\ref{thm:transfer}).

Huang et al.~\cite{reasoning_jailbreak2026} demonstrate that large reasoning models
achieve a $97\%$ attack success rate when deployed as autonomous jailbreak agents,
confirming that LLM reasoning capability itself is the primary attack surface.
This directly parallels our core thesis: the attack does not exploit model
miscalibration but the model's genuine reasoning competence applied to
abstractions the guardrail cannot interpret.

\subsection{Encoding and Structural Attacks}\label{sec:related-encoding}

StructuralSleight~\cite{structuralsleight2024} achieves $94.62\%$ attack success
on GPT-4o via uncommon text-encoded structures (ASCII art, cipher text, etc.),
demonstrating that encoding diversity defeats content filters.
Our algebraic results explain \emph{why}: any encoding that requires modular
counting to decode lies outside the $\mathrm{AC}^0$ closure of any aperiodic
guardrail language (Theorem~\ref{thm:blindness}).

Boucher et al.~\cite{boucher2021badchars} demonstrate imperceptible attacks via
Unicode manipulation---inserting invisible characters between payload characters
so that the guardrail sees an interrupted substring.
Zero-width character insertion is an instance of $\mathrm{MOD}_2$ interleaving with
a null filler character; our Corollary~\ref{cor:parity} proves it cannot be patched
within the regex paradigm.

The Content Concretization attack~\cite{concretization2025} traverses the
abstraction hierarchy in the opposite direction: given an abstract harmful
specification, the attacker prompts the LLM to supply concrete details.
We proceed concrete-to-abstract---encoding a concrete harmful task as an abstract
formal system---and the interpretation happens offline.
The two approaches are complementary instances of the same general principle that
abstraction boundaries defeat content inspection.

DictAttack~\cite{dictattack2025} embeds malicious intent in grammar-level schema rules
while maintaining benign surface prompts, achieving 94--99\% attack success rates across
multiple LLMs.
This structured-output attack exploits the same abstraction boundary our V3 and V4 vectors
formalize: the gap between surface-level content filtering and structural information
encoding.
Our contribution is the formal impossibility proof that regex guardrails \emph{cannot}
close this gap, whereas DictAttack provides empirical evidence of its exploitability.

S2C~\cite{s2c2026} is independent concurrent work (March 2026) that manipulates
semantic reconstruction during inference across multiple semantic dimensions,
arriving at a complementary abstraction-based framing through a different research
path.
Where S2C operates within a single LLM inference pass, our pipeline distributes
the harmful reasoning across multiple agents and an offline interpretation table;
both works independently identify abstraction boundaries as the core attack surface.

\subsection{Computational and Algebraic Impossibility}\label{sec:related-impossibility}

Ball et al.~\cite{ball2025impossibility} prove that under standard cryptographic
assumptions, adversarial prompts can be computationally indistinguishable from
benign ones for any efficient filter.
Our results are complementary.
We establish \emph{syntactic} (perfect, unconditional) indistinguishability for
formal reasoning attacks (Proposition~\ref{thm:indistinguish}) and \emph{algebraic}
(unconditional) blindness for modular counting attacks
(Theorem~\ref{thm:blindness}).
Ball et al.\ give a more general but assumption-dependent result over a broader
attack class; our results apply to specific attack constructions and carry no
computational hardness assumption.
Kumar et al.~\cite{certifying_safety2023} propose certified defenses via randomized
smoothing applied to LLM inputs, establishing provable safety radii around benign
inputs under $\ell_p$-bounded perturbations; our results identify a complementary
regime where the attack is not a small perturbation but a structure-preserving map
that lies outside the support of the smoothing distribution.
Dong et al.~\cite{nfl_guardrails2025} demonstrate a no-free-lunch tradeoff between
guardrail thoroughness and computational cost, showing that any guardrail imposes
overhead that scales with threat coverage; our results show that even unbounded
overhead cannot patch the regex blindness established in Theorem~\ref{thm:blindness}.
R2-Guard~\cite{r2guard2025} replaces regex matching with reasoning-based evaluation
via chain-of-thought; this approach operates in $\mathrm{TC}^0$ and is not subject
to the aperiodic blindness theorem, but remains susceptible to the information-theoretic
and structural indistinguishability barriers established in
\S\S\ref{sec:fano}--\ref{sec:indistinguish}.

The abstract encoding channel in homomorphic reasoning is a form of
\emph{steganographic channel}~\cite{hopper2002steganography}: the payload
(harmful intent) is hidden in a covertext (legitimate formal reasoning task)
such that no efficient observer (the guardrail) can distinguish stegotext from
covertext.
Hopper, Langford, and von Ahn~\cite{hopper2002steganography} formalize perfect
steganographic security as the requirement that no polynomial-time distinguisher
achieves better than negligible advantage in separating channel transmissions from
covertext.
Proposition~\ref{thm:indistinguish} can be viewed as establishing \emph{perfect}
steganographic security for the homomorphic reasoning channel: the distinguisher
advantage is exactly zero (not merely negligible) because the LLM-visible content
of the attack is syntactically equal to a legitimate reasoning task, making the
distinguisher's task information-theoretically impossible rather than merely
computationally hard.

\subsection{Formal Languages and Neural Networks}\label{sec:related-formal}

Merrill~\cite{merrill2024survey} surveys the formal language expressivity of
transformer architectures.
Merrill and Sabharwal~\cite{merrill2023saturated} establish that log-precision
transformers are contained in $\mathrm{TC}^0$, implying they can compute modular
counting functions.
Chiang et al.~\cite{chiang2024tighter} give tighter expressivity bounds.
Hahn~\cite{hahn2020limitations} proves that self-attention mechanisms are limited
on periodic formal languages, providing circuit-theoretic evidence for the same
qualitative boundary our algebraic results formalize.
Weiss et al.~\cite{weiss2022extracting} extract DFAs from RNNs using the Angluin
$L^*$ algorithm~\cite{angluin1987}; the same DFA extraction pipeline underlies our
empirical syntactic monoid computation.

\subsection{Indirect Prompt Injection}\label{sec:related-injection}

Greshake et al.~\cite{greshake2023injection} establish that LLM-integrated
applications blur the boundary between data and instructions, enabling an attacker
to inject control flow via data channels.
Our V2 (zero-knowledge pipeline) is a specific instance of this pattern: the
``data'' is a base64-encoded configuration that contains harmful parameters the
agent processing it never reads.
The theoretical underpinning differs---Greshake et al.\ exploit instruction/data
confusion, while V2 exploits the information destruction formalized in
Proposition~\ref{prop:fano}---but the practical attack pattern is closely related.

\subsection{Tool-Chain Composition Attacks}\label{sec:related-toolchain}

Zhou et al.~\cite{stac2025} formalize the STAC (Simple Tool-Augmented Chain)
attack class, in which individually benign tool calls compose into harmful
pipelines---achieving $>90\%$ attack success rate across 483 attack cases.
Our V1 (decomposition) and V2 (zero-knowledge pipeline) are specific instances
of this general pattern.
The distinction is that STAC provides a large-scale empirical evaluation of
tool-chain composition, while our contribution is the formal proof
(Proposition~\ref{prop:fano}, Proposition~\ref{thm:np}) that no inference-layer
guardrail can detect such compositions without inspecting runtime state.

% ============================================================
\section{Threat Model and Formalization}\label{sec:threat}

\subsection{Attacker and Defender Model}\label{sec:threat-model}

\paragraph{Attacker capabilities.}
The attacker has access to:
(a)~an LLM-based coding assistant with bash execution (\eg Claude Code, Copilot, Cursor);
(b)~local LLM infrastructure (\eg Ollama with open-weight models) for offline
grammar solving without guardrail exposure;
(c)~standard programming tools (Python, AST manipulation, base64 encoding).

\paragraph{Defender capabilities.}
The guardrail system operates at LLM inference time via RLHF training,
system prompt rules, and output filters.
It inspects the content visible to the LLM---prompts, generated code, tool
calls---but not runtime arguments, environment variables, base64-encoded data
in files the agent does not read, or offline interpretation tables.

\paragraph{Attacker goal.}
Cause the LLM to produce computation that, when instantiated with
attacker-controlled parameters or interpreted via an attacker-held interpretation
table, performs a task the guardrail would refuse if asked directly.
The LLM must never observe the harmful intent in plaintext during any inference call
subject to guardrail inspection.

\paragraph{Out of scope.}
Attacks requiring model weight access, fine-tuning, or gradient computation.
Attacks against execution-layer defenses (sandboxing, network ACLs).
Social engineering of human reviewers.

\subsection{Abstraction-Based Formalization}\label{sec:formalization}

\begin{definition}[Abstraction function]\label{def:abstraction}
Let $R$ be a finite set of concrete artifacts and $M$ a finite set of abstract types.
An \emph{abstraction function} $U\colon R \to M$ maps each concrete artifact to its
abstract type (\eg $U(\texttt{streaming.example}) = \text{``video endpoint''}$,
$U(\texttt{C}_9\texttt{H}_8\texttt{O}_4) = \text{``acetyl compound''}$).
\end{definition}

We do not require $U$ to be surjective; we only use the existence of the fiber
structure defined below.

\begin{definition}[Fiber]\label{def:fiber}
For $m \in M$, the \emph{fiber} $U^{-1}(m) = \{r \in R : U(r) = m\}$ is the
set of all concrete artifacts mapping to abstract type $m$.
A fiber is \emph{mixed} if it contains both allowed ($G(r) = 0$) and denied
($G(r) = 1$) artifacts under the ground-truth guardrail $G\colon R \to \{0,1\}$.
\end{definition}

\begin{definition}[Lifted guardrail]\label{def:lifted}
A \emph{lifted guardrail} is a function $G'\colon M \to \{0,1\}$ that approximates
the ground-truth guardrail $G$ through the abstraction layer: the system classifies
concrete artifact $r$ by evaluating $G'(U(r))$ rather than $G(r)$ directly.
\end{definition}

The existence of mixed fibers is the root cause of the information-theoretic lower
bound in Proposition~\ref{prop:fano}: if a fiber $U^{-1}(m)$ contains both safe
and unsafe concrete artifacts, no lifted guardrail $G'$ can correctly classify all
of them without access to $r$ itself---but the guardrail operating at the
abstraction layer sees only $m = U(r)$.

\subsection{Guardrails as Language Recognizers}\label{sec:guardrail-language}

Let $\Sigma$ be the finite alphabet of LLM prompt characters.

\begin{definition}[Regex guardrail]\label{def:regex-guardrail}
A \emph{regex guardrail} is a function $G\colon \Sigma^* \to \{\mathsf{allow}, \mathsf{block}\}$
defined by a finite set of regular expression patterns $\{r_1, \ldots, r_k\}$
applied via substring search:
\[
  G(w) = \begin{cases}
    \mathsf{block} & \text{if } \exists\, i,\; \exists\, u,v \in \Sigma^*:
                      w = u \cdot x \cdot v,\; x \in L(r_i), \\
    \mathsf{allow} & \text{otherwise.}
  \end{cases}
\]
\end{definition}

Let $L_{\mathsf{safe}} = \{w \in \Sigma^* : G(w) = \mathsf{allow}\}$ be the set of
strings the guardrail permits.

\begin{definition}[Blocked language]\label{def:blocked-language}
The \emph{blocked language} of $G$ is the complement
$L_G = \Sigma^* \setminus L_{\mathsf{safe}} = \Sigma^* \cdot \bigl(\bigcup_{i=1}^{k} L(r_i)\bigr) \cdot \Sigma^*$.
This is a regular language (closed under concatenation with $\Sigma^*$ and finite union).
\end{definition}

\begin{definition}[Guardrail blindness to encoding class]\label{def:blindness}
Let $\mathcal{E}$ be a class of invertible encodings $\phi\colon \Sigma^* \to \Sigma^*$.
A guardrail $G$ is \emph{blind to encoding class $\mathcal{E}$} if for every
$\phi \in \mathcal{E}$:
\begin{enumerate}
  \item $\phi(L_{\mathsf{blocked}}) \cap L_G = \emptyset$ (encoded harmful strings
    pass the guardrail), and
  \item there is no regular language $L'$ with aperiodic syntactic monoid such that
    $\phi(L_{\mathsf{blocked}}) \subseteq L' \subseteq \overline{L_{\mathsf{safe}}}$
    (no aperiodic regex can separate encoded-harmful from genuinely-benign strings).
\end{enumerate}
\end{definition}

Condition (1) states that the encoding evades the current guardrail; condition (2)
states that the evasion cannot be patched by adding any regex whose syntactic monoid
is aperiodic---which by \S\ref{sec:algebraic-foundations} is the entire aperiodic
regex class.
This pair of conditions connects the regex language-recognizer formalization to the
abstraction-based formalization of \S\ref{sec:formalization}: regex guardrails are
a specific instance of lifted guardrails operating on string content, and their
fibers are the equivalence classes of the syntactic congruence $\equiv_{L_G}$.

% ============================================================
\section{Algebraic Foundations}\label{sec:algebraic-foundations}

This section recalls the classical algebraic and circuit-theoretic machinery
underlying the impossibility results of \S\ref{sec:blindness-theorem}.
All results here are classical; full developments appear in
Pin~\cite{pin1986}, Eilenberg~\cite{eilenberg1976}, and
Straubing~\cite{straubing1994}.

\subsection{Syntactic Monoids}\label{sec:syntactic-monoids}

For a finite alphabet $\Sigma$ and regular language $L \subseteq \Sigma^*$,
the \emph{syntactic congruence} $\equiv_L$ identifies two strings when they
behave identically in every left-right context: $u \equiv_L v$ iff
$xuy \in L \Leftrightarrow xvy \in L$ for all $x, y \in \Sigma^*$.

\begin{definition}[Syntactic monoid]\label{def:syntactic-monoid}
The \emph{syntactic monoid} of $L$ is $M(L) = \Sigma^* / {\equiv_L}$, with the
operation inherited from string concatenation.
It is finite~\cite{myhill1957,nerode1958}, isomorphic to the transition monoid
of the minimal DFA for $L$, and the smallest monoid recognizing $L$ up to
division~\cite{pin1986,eilenberg1976}---hence computable from the minimal DFA
by standard monoid enumeration.
\end{definition}

\subsection{Aperiodicity and Star-Freeness}\label{sec:aperiodicity}

\begin{definition}[Aperiodic monoid]\label{def:aperiodic}
A monoid $M$ is \emph{aperiodic} (equivalently, \emph{group-free}) if it
contains no nontrivial subgroup: for all $m \in M$, there exists $n \geq 1$
with $m^n = m^{n+1}$.
Cyclic groups $\mathbb{Z}/p\mathbb{Z}$ are the prototypical non-aperiodic
monoids; aperiodic monoids cannot replicate their cyclic behavior.
\end{definition}

\begin{theorem}[Sch\"{u}tzenberger 1965, McNaughton--Papert 1971~{\cite{schutzenberger1965,mcnaughton1971}}]\label{thm:smp}
For a regular language $L$, the following are equivalent:
(i)~$M(L)$ is aperiodic;
(ii)~$L$ is \emph{star-free}---definable from $\Sigma$, $\emptyset$,
$\{\varepsilon\}$, and single letters using Boolean operations and
concatenation, without Kleene star;
(iii)~$L$ is definable in first-order logic with linear order,
$\mathrm{FO}[{<}]$.
\end{theorem}

Star-freeness and $\mathrm{FO}[{<}]$-definability are ``counting-free'':
they express positional and ordering constraints but not modular periodicity.

\subsection{Circuit Complexity and Modular Counting}\label{sec:circuit-complexity}

$\mathrm{AC}^0$ is the class of languages recognized by families of
constant-depth, polynomial-size Boolean circuits with unbounded fan-in AND/OR
gates---the computations that parallelize without counting.

\begin{theorem}[Barrington--Compton--Straubing--Th\'{e}rien 1992~{\cite{barrington1992}}]\label{thm:bcst}
A regular language is in $\mathrm{AC}^0$ iff its syntactic monoid is aperiodic.
\end{theorem}

\begin{theorem}[Furst--Saxe--Sipser 1981, H\r{a}stad 1987~{\cite{furst1981,hastad1987}}]\label{thm:fss}
$\mathrm{PARITY} \notin \mathrm{AC}^0$; more generally, $\mathrm{MOD}_p \notin
\mathrm{AC}^0$ for any prime $p$, where $\mathrm{MOD}_p = \{w \in \{0,1\}^* :
|w|_1 \equiv 0 \pmod{p}\}$.
\end{theorem}

Chaining Theorems~\ref{thm:smp}--\ref{thm:fss}: aperiodic $M(L_G) \Rightarrow
L_G \in \mathrm{AC}^0 \Rightarrow L_G$ cannot compute $\mathrm{MOD}_p$, making
the guardrail blind to any payload encoding whose decoder requires modular
counting.

\subsection{Krohn--Rhodes Decomposition}\label{sec:krohn-rhodes}

\begin{theorem}[Krohn--Rhodes 1965~{\cite{krohnrhodes1965}}]\label{thm:kr-classical}
Every finite monoid divides a wreath product of simple groups and aperiodic
semigroups; the simple groups appearing in the decomposition (the
\emph{group components}) are uniquely determined up to isomorphism.
\end{theorem}

The group components of $M(L_G)$ form a \emph{complete spectrum of algebraic
blind spots}: by the variety-theoretic characterization of
Eilenberg~\cite{eilenberg1976} and Straubing~\cite{straubing1994}, $L_G$ can
potentially express $\mathrm{MOD}_p$ only if some group component has order
divisible by $p$; the guardrail is provably blind to all other primes.
When $M(L_G)$ is aperiodic (no group components) blindness is total---the
case for every substring-matching regex guardrail, as established by
Theorem~\ref{thm:substring-aperiodic}.

% ============================================================
\section{Impossibility Results}\label{sec:blindness-theorem}

\subsection{Substring-Matching Patterns Are Aperiodic}\label{sec:substring-aperiodic}

\begin{theorem}[Substring aperiodicity]\label{thm:substring-aperiodic}
Let $r$ be a regular expression built from literals, character classes,
bounded repetition, alternation, optional groups, dot, and Kleene star or plus
applied to character classes.
Let $L_r = \Sigma^* \cdot L(r) \cdot \Sigma^*$ be the language of strings
containing a substring matching $r$.
Then the syntactic monoid $M(L_r)$ is aperiodic.
\end{theorem}

\begin{proof}
We show $L_r$ is star-free, hence aperiodic by the equivalence in
Theorem~\ref{thm:smp}.

\paragraph{Star-free closure under substring wrapping.}
For any regular language $K$ over $\Sigma$, the language
$\Sigma^* \cdot K \cdot \Sigma^*$ is star-free whenever $K$ is star-free.
This follows directly from the closure of star-free languages under
concatenation~\cite{pin1986}: $\Sigma^* = \overline{\emptyset}$ is star-free,
and the concatenation of three star-free languages is star-free.
It therefore suffices to show that $L(r)$ itself is star-free for each
building-block construction.

\paragraph{Literals and character classes.}
A single character $a \in \Sigma$ defines $\{a\}$, trivially star-free.
A character class $[c_1 c_2 \cdots c_k]$ is $\{c_1\} \cup \cdots \cup \{c_k\}$,
a finite union of singletons, hence star-free.

\paragraph{Bounded repetition.}
$r^{\{n\}}$ (exactly $n$ repetitions) is the $n$-fold concatenation of $L(r)$,
star-free by induction on $n$ using closure under concatenation.
$r^{\{m,n\}}$ (between $m$ and $n$ repetitions) is $\bigcup_{k=m}^{n} r^{\{k\}}$,
a finite union of star-free languages, hence star-free.

\paragraph{Alternation.}
$r_1 \mid r_2$ yields $L(r_1) \cup L(r_2)$; star-free languages are closed under
union (and Boolean operations generally), so this is star-free.

\paragraph{Optional groups.}
$r?$ yields $L(r) \cup \{\varepsilon\}$; since $\{\varepsilon\}$ is star-free and
star-free is closed under union, this is star-free.

\paragraph{Dot.}
The pattern $.$ matches any single character in $\Sigma$.
It equals $\bigcup_{a \in \Sigma}\{a\}$, a finite union of singletons, star-free.

\paragraph{Kleene plus over character classes.}
Let $C$ be a character class with language $L(C) \subseteq \Sigma$.
We claim $L(C)^+$ is star-free.
$L(C)^+$ is the set of nonempty strings over $L(C)$, equal to
$\overline{\overline{L(C) \cdot \Sigma^*} \cup \{\varepsilon\}}$; more directly,
$L(C)^+ = \Sigma^* \setminus (\Sigma^* \cdot \overline{L(C)} \cdot \Sigma^* \cup \{\varepsilon\})$.
Since $\overline{L(C)} = \Sigma \setminus L(C)$ is a finite set of singletons (hence
star-free), and star-free is closed under complement and concatenation with $\Sigma^*$,
$L(C)^+$ is star-free.
The argument applies equally to $L(C)^*$: $L(C)^* = L(C)^+ \cup \{\varepsilon\}$,
a union of two star-free languages.

\paragraph{Conclusion.}
Every building block of $r$ produces a star-free language, and
closure under concatenation and Boolean operations preserves star-freeness
throughout the construction of $L(r)$.
Wrapping with $\Sigma^* \cdot (\cdot) \cdot \Sigma^*$ preserves star-freeness.
Therefore $L_r$ is star-free, and by Theorem~\ref{thm:smp} its syntactic
monoid $M(L_r)$ is aperiodic.
\end{proof}

\begin{remark}[Guardrail patterns vs.\ $(aa)^*$]
A pattern such as $(aa)^*$ recognizes strings of even $a$-count, requiring
$\mathbb{Z}/2\mathbb{Z}$ in its syntactic monoid, so it is \emph{not} aperiodic.
However, $(aa)^*$ is not a substring-matching guardrail pattern: it is a
character-class Kleene iteration over a two-character literal, not a single
character class.
The guardrail patterns in our corpus---keyword blocklists, domain
patterns, toxicity keywords---are constructed precisely as described in the
theorem: Kleene star and plus are applied to character classes, not to
multi-character sequences.
The syntactic monoid of the representative guardrail pattern $(ab)^*$ is the
six-element monoid $B_2 = \{1, a, b, ab, ba, 0\}$ where multiplication follows
from string concatenation modulo the syntactic congruence; this monoid has no
nontrivial subgroups, confirmed by our monoid extractor (Section~\ref{sec:corpus}).
\end{remark}

\subsection{The Guardrail Blindness Theorem}\label{sec:blindness-proof}

Define the \emph{modular decode language} as follows.
Fix a prime $p \geq 2$.
Let $\Sigma_0 = \{0, 1\}$ and let $w \in \{0,1\}^*$.
The \emph{$p$-interleaved encoding} of a harmful string $h \in L_{\mathsf{blocked}}$
is a string $\phi_p(h) \in \Sigma^*$ constructed by inserting $p-1$ arbitrary
filler characters between successive characters of $h$, so that position $i$
(zero-indexed) of $h$ appears at position $p \cdot i$ in $\phi_p(h)$.
The decoder recovers $h$ by reading every $p$-th character: $h = \phi_p(h)[0::p]$.
Define
\[
  L_{\mathsf{decode}}(p)
  = \bigl\{ w \in \Sigma^* :
    w[0::p] \in L_{\mathsf{blocked}} \bigr\}
\]
as the language of strings whose every-$p$-th character subsequence is in the
blocked language.

\begin{theorem}[Guardrail blindness]\label{thm:blindness}
Let $G$ be any regex guardrail with aperiodic syntactic monoid $M(L_G)$.
Let $p \geq 2$ be any prime.
Then:
\begin{enumerate}
  \item $L_{\mathsf{decode}}(p)$ is regular.
  \item $M(L_{\mathsf{decode}}(p))$ is not aperiodic; it contains
        $\mathbb{Z}/p\mathbb{Z}$ as a subgroup.
  \item $\phi_p(L_{\mathsf{blocked}}) \cap L_G = \emptyset$, and no
        modification of $G$ within the class of aperiodic regex guardrails
        can detect $L_{\mathsf{decode}}(p)$.
\end{enumerate}
\end{theorem}

\begin{proof}\leavevmode
\paragraph{(1) Regularity.}
The decoding function $w \mapsto w[0::p]$ is computed by a length-$p$ cyclic DFA
that counts position modulo $p$ and records the character at position $\equiv 0
\pmod{p}$.
This DFA has state set $Q = \{0, 1, \ldots, p-1\} \times Q_{\mathsf{block}}$,
where $Q_{\mathsf{block}}$ is the state set of the minimal DFA for $L_{\mathsf{blocked}}$.
Transitions increment the counter modulo $p$ and advance the $L_{\mathsf{blocked}}$
component only when the counter is $0$.
Since both components are finite, $L_{\mathsf{decode}}(p)$ is accepted by a finite
automaton of size $p \cdot |Q_{\mathsf{block}}|$ and is therefore regular.

\paragraph{(2) Non-aperiodicity.}
Consider the DFA constructed above, with states $(i, q) \in \{0,\ldots,p{-}1\} \times Q_{\mathsf{block}}$.
Let $f \in \Sigma$ be a \emph{filler} character that appears only at non-zero counter
positions---i.e., $\delta_{\mathsf{block}}(q, f) = q$ for all $q \in Q_{\mathsf{block}}$
(the $Q_{\mathsf{block}}$ component is unchanged by $f$).
The transition induced by reading $f$ maps $(i, q) \mapsto ((i+1) \bmod p,\, q)$,
advancing only the counter.
Consider the word $w = f^p$ consisting of $p$ copies of $f$.
For any fixed $q \in Q_{\mathsf{block}}$, the transformation $\delta_w$ maps
$(i, q) \mapsto ((i + p) \bmod p,\, q) = (i, q)$; thus $\delta_w$ acts as the
identity on states of the form $(i, q)$ for this $q$.
However, $\delta_{f^j}$ maps $(i, q) \mapsto ((i+j) \bmod p,\, q)$ for $1 \le j < p$,
which is not the identity since $j \not\equiv 0 \pmod{p}$.
Therefore $\delta_{f^j} \neq \mathrm{id}$ for $1 \le j < p$, yet $\delta_{f^p} = \mathrm{id}$
on every state of the form $(i, q)$ with $q$ fixed.
The elements $\{\delta_{f^j} : 0 \le j < p\}$ form a subgroup of the transition
monoid isomorphic to $\mathbb{Z}/p\mathbb{Z}$, witnessed by the cyclic permutation
$(i, q) \mapsto ((i+1) \bmod p,\, q)$ of order exactly $p$.
The cyclic group $\mathbb{Z}/p\mathbb{Z}$ therefore divides $M(L_{\mathsf{decode}}(p))$,
which is thus non-aperiodic.

\paragraph{(3) Blindness consequence.}
By construction, $\phi_p(L_{\mathsf{blocked}}) \subseteq L_{\mathsf{decode}}(p)$.
If $w = \phi_p(h)$ for some $h \in L_{\mathsf{blocked}}$, then $w$ contains no
substring in $\bigcup_i L(r_i)$ (the harmful keywords appear only at every $p$-th
position, not as contiguous substrings), so $G(w) = \mathsf{allow}$.

For the patchability claim: any regex guardrail added to $G$ with an aperiodic
syntactic monoid defines a language in $\mathrm{AC}^0$ by Theorem~\ref{thm:bcst}.
But detecting whether $w[0::p] \in L_{\mathsf{blocked}}$ requires computing a
$\mathrm{MOD}_p$ predicate on the position counter, which is not in $\mathrm{AC}^0$
by Theorem~\ref{thm:fss} (for prime $p$).
Therefore, no aperiodic regex can be added that correctly classifies all encodings
in $\phi_p(L_{\mathsf{blocked}})$ while leaving genuinely benign strings allowed.
\end{proof}

\begin{corollary}[$\mathrm{MOD}_2$ sufficiency]\label{cor:parity}
$p = 2$ suffices for a complete bypass: every aperiodic regex guardrail is blind to
the alternating-character encoding $\phi_2$, which interleaves the payload with
arbitrary filler characters at every odd position.
Existing evasion techniques---reading every other character, zero-width character
insertion between payload characters, even-position acrostics---are all instances
of $\phi_2$ and are therefore provably undetectable by any current regex guardrail.
\end{corollary}

\begin{remark}[Extension to composite moduli]
Theorem~\ref{thm:blindness} is stated for prime $p$.
For composite $p = q_1 \cdots q_k$, blindness to $\mathrm{MOD}_p$ follows from
blindness to each prime factor: if the guardrail cannot detect
$\mathrm{MOD}_{q_i}$ for any $i$, it cannot detect $\mathrm{MOD}_p$.
The prime case therefore implies the general case for aperiodic guardrails,
which are blind to all primes simultaneously.
\end{remark}

\subsection{Composition Preserves Aperiodicity}\label{sec:composition}

\begin{proposition}[Composition preserves aperiodicity]\label{prop:composition}
If $L_1, L_2 \subseteq \Sigma^*$ are regular languages with aperiodic syntactic monoids,
then $L_1 \cap L_2$, $L_1 \cup L_2$, and $\Sigma^* \setminus L_1$ all have aperiodic syntactic monoids.
\end{proposition}
\begin{proof}[Proof sketch]
The syntactic monoid of $L_1 \cap L_2$ divides $M(L_1) \times M(L_2)$.
A divisor of a direct product of aperiodic monoids is aperiodic,
since aperiodicity is preserved under division (a divisor of an aperiodic monoid has no nontrivial subgroups).
The cases of union and complement follow from De Morgan's laws and the fact that $M(\overline{L}) \cong M(L)$.
\end{proof}

This means that applying multiple guardrail patterns in sequence (intersection) or
in parallel (union) preserves the blindness property.
A guardrail system composed of aperiodic components remains blind to all modular encodings.

\subsection{Krohn--Rhodes Attack Taxonomy}\label{sec:kr-taxonomy}

\begin{proposition}[Blindness spectrum, after Straubing~\cite{straubing1994}]\label{prop:krohn-rhodes}
Let $G$ be a regex guardrail with syntactic monoid $M(L_G)$ and Krohn--Rhodes
group components $G_1, G_2, \ldots, G_k$ (simple groups appearing in the
decomposition of Theorem~\ref{thm:kr-classical}).
Then:
\begin{enumerate}
  \item $G$ \emph{can detect} $\mathrm{MOD}_p$-encoded payloads only if
        $p$ divides $|G_i|$ for some $i \in \{1,\ldots,k\}$.
  \item $G$ is \emph{provably blind} to $\mathrm{MOD}_q$-encoded payloads for
        every prime $q$ not dividing any $|G_i|$.
  \item If $M(L_G)$ is aperiodic (no group components), then $G$ is blind to
        $\mathrm{MOD}_p$-encoded payloads for every prime $p \geq 2$---blindness
        is total.
\end{enumerate}
\end{proposition}

\begin{remark}[Constructive audit via monoid extractor]
The Krohn--Rhodes blindness spectrum is constructively computable.
Given a guardrail regex, the monoid extractor (Section~\ref{sec:corpus}) computes
the minimal DFA, enumerates the transition monoid, and identifies all group
components via standard semigroup algorithms.
The output is a complete list of primes to which the guardrail is blind: if the
group components are $\{S_5, \mathbb{Z}/3\mathbb{Z}\}$, the guardrail can
potentially detect $\mathrm{MOD}_3$ and $\mathrm{MOD}_5$ encodings but is blind to
all other modular encodings.
For $100\%$ of patterns in our corpus, the component list is empty and blindness
is total.
\end{remark}

\subsection{An Illustrative Lower Bound (Fano)}\label{sec:fano}

We include an information-theoretic lower bound as a secondary, illustrative
result---presented under a simplifying uniformity assumption, and intended to
motivate rather than to match the force of the algebraic blindness theorems of
\S\ref{sec:blindness-proof}.
The bound is of independent interest but should not be read as an unconditional
impossibility: a distribution-aware version would require modeling the
attacker's prompt distribution (see \S\ref{sec:limitations-list}, item~3).

\begin{proposition}[Information destruction]\label{prop:fano}
Let $U\colon R \to M$ be an abstraction function and $G\colon R \to \{0,1\}$
the ground-truth guardrail.
Let $R$ be distributed uniformly over $R$.
Then for any lifted guardrail $G'\colon M \to \{0,1\}$, the probability of
misclassification satisfies
\[
  P_e \;=\; \Pr[G'(U(R)) \neq G(R)]
  \;\geq\;
  h^{-1}\!\bigl(H(G(R) \mid U(R))\bigr),
\]
where $h(x) = -x\log_2 x - (1-x)\log_2(1-x)$ is the binary entropy function and
$h^{-1}$ is its inverse on $[0, \tfrac{1}{2}]$.
\end{proposition}

\begin{proof}
Consider the Markov chain $G(R) \to U(R) \to G'(U(R))$.
By the data-processing inequality, the mutual information satisfies
$I(G(R); G'(U(R))) \leq I(G(R); U(R)) = H(G(R)) - H(G(R) \mid U(R))$.
The lifted guardrail's decision $G'(U(R))$ is a binary function of $U(R)$, so the
probability of error is $P_e = \Pr[G'(U(R)) \neq G(R)]$.
By Fano's inequality~\cite{cover2006elements},
$H(G(R) \mid G'(U(R))) \leq h(P_e) + P_e \cdot \log_2(|\{0,1\}| - 1) = h(P_e)$.
Since conditioning can only reduce entropy,
$H(G(R) \mid U(R)) \leq H(G(R) \mid G'(U(R))) \leq h(P_e)$.
Inverting the (monotone increasing on $[0, \tfrac{1}{2}]$) binary entropy function
gives $P_e \geq h^{-1}(H(G(R) \mid U(R)))$.
\end{proof}

\paragraph{Interpretation.}
The conditional entropy $H(G(R) \mid U(R))$ is positive precisely when mixed fibers
exist (Definition~\ref{def:fiber}).
As an illustrative calculation under the uniformity assumption: if $30\%$ of
probability mass falls in mixed fibers with ${\sim}50/50$ safe/unsafe split, then
$H(G(R) \mid U(R)) \approx 0.3$ bits, yielding
$P_e \ge h^{-1}(0.3) \approx 0.048$ (about $5\%$).
Under uniformity the floor is a genuine information-theoretic limit, but the
quantitative bound is sensitive to the prior: real prompt distributions are
highly non-uniform, and the floor may tighten or loosen substantially depending
on how adversarial mass is distributed across fibers (see
\S\ref{sec:limitations-list}, item~3).
The qualitative message---reducing abstraction to restore fidelity makes the
guardrail brittle to novel formulations---is robust to the distributional
caveat, but the numerical floor should be read illustratively.

\subsection{Computational Intractability (NP-Completeness)}\label{sec:np}

\begin{definition}[Guardrail Instantiation Decision Problem (GIDP)]\label{def:gidp}
Given:
\begin{itemize}
  \item An abstract program $P$ over a finite set of abstract symbols $M$,
  \item An interpretation domain $R$ with ground-truth guardrail $G\colon R \to \{0,1\}$,
  \item An interpretation function schema $I\colon M \to 2^R$ mapping each abstract
    symbol to a nonempty set of candidate concrete values,
\end{itemize}
\textsc{GIDP} asks: does there exist a \emph{total assignment}
$\iota\colon M \to R$ with $\iota(m) \in I(m)$ for all $m$ such that the concretized
program $P[\iota]$ satisfies $G(P[\iota]) = 1$ (i.e., the program is dangerous)?
\end{definition}

\begin{proposition}[NP-completeness of GIDP]\label{thm:np}
\textsc{GIDP} is NP-complete.
This follows from a standard reduction from 3-SAT.
\end{proposition}

\begin{proof}\leavevmode
\paragraph{Membership in NP.}
A certificate is an assignment $\iota\colon M \to R$.
Given $\iota$, checking $\iota(m) \in I(m)$ for all $m$ takes time $O(|M| \cdot |R|)$;
evaluating $G(P[\iota])$ takes polynomial time in the program size.
Therefore \textsc{GIDP} $\in$ NP.

\paragraph{NP-hardness.}
We reduce 3-SAT to \textsc{GIDP}.
Given a 3-CNF formula $\varphi$ with variables $x_1, \ldots, x_n$ and clauses
$C_1, \ldots, C_m$, construct:
\begin{itemize}
  \item Abstract symbols $M = \{v_1, \ldots, v_n\}$, one per variable.
  \item Concrete domain $R = \{T, F\}$ (Boolean values).
  \item $I(v_i) = \{T, F\}$ for each $i$: each variable can be assigned
    either truth value.
  \item A total assignment $\iota\colon M \to R$ with $\iota(v_i) \in I(v_i)$
    corresponds to a Boolean assignment $x_i = \iota(v_i)$.
  \item Abstract program $P = (v_1, \ldots, v_n)$ is the ordered list of
    variable symbols.
  \item Guardrail $G\colon R^n \to \{0,1\}$ is the Boolean circuit encoding
    $\varphi$: $G(\iota(v_1), \ldots, \iota(v_n)) = 1$ iff
    $\varphi(\iota(v_1), \ldots, \iota(v_n))$ is satisfiable.
\end{itemize}
A satisfying assignment for $\varphi$ exists iff the \textsc{GIDP} instance has
a positive answer.
Since 3-SAT is NP-complete, \textsc{GIDP} is NP-hard.
\end{proof}

\begin{remark}[Pointwise vs.\ combinatorial; scope of hardness]
The NP-hardness comes from the combinatorial search over the exponentially large
space of concrete interpretations, not from any single concretization step.
When $G$ is restricted to the class of regex guardrails (Definition~\ref{def:regex-guardrail}),
the single-interpretation check $G(P[\iota])$ is computable in time
$O(|P| \cdot |Q_G|)$ where $|Q_G|$ is the DFA state count.
The NP-hardness therefore resides entirely in the search over the product of fiber
sets $|I(m_1)| \times \cdots \times |I(m_k)|$, not in the evaluation of the guardrail
itself.
If the interpretation function $I$ is a bijection (each abstract symbol has a
unique concrete value), \textsc{GIDP} is solvable in polynomial time: just evaluate
$G(P[\iota])$ directly.
The hardness is precisely the problem of searching through a product of fiber
sets---the same mixed-fiber phenomenon that drives the Fano bound.
\end{remark}

\subsection{Faithful Transfer (Free Category Functoriality)}\label{sec:functor}

\begin{proposition}[Faithful transfer]\label{thm:transfer}
Let $\mathcal{A}$ be the free category generated by a context-free grammar
$\mathcal{G}$ over abstract symbols $M$, and let $\mathcal{D}$ be the category of
domain operations over $R$.
Let $I\colon M \to R$ be a graph homomorphism from the grammar's dependency
structure to $\mathcal{D}$.
Then $I$ extends uniquely to a functor $\hat{I}\colon \mathcal{A} \to \mathcal{D}$
that maps every derivation path in $\mathcal{A}$ to a valid operation sequence in
$\mathcal{D}$.
If $I$ is additionally injective on the generating morphisms of $\mathcal{G}$, then
$\hat{I}$ is faithful: distinct derivation paths map to distinct operation sequences.
This is a direct application of the universal property of free categories \cite{maclane1971categories}.
\end{proposition}

\begin{proof}
$\mathcal{A}$ is the free category on the graph $\mathcal{G}$.
By the universal property of free categories~\cite{maclane1971categories}
(Section~II.7), any graph homomorphism $I$ into a category $\mathcal{D}$ extends
uniquely to a functor $\hat{I}\colon \mathcal{A} \to \mathcal{D}$.
Functoriality ensures composition preservation: multi-step derivations map to
multi-step domain operations.
For faithfulness: if $I$ is injective on generating morphisms, then distinct
generators $f_{m_1} \neq f_{m_2}$ map to distinct operations
$\hat{I}(f_{m_1}) \neq \hat{I}(f_{m_2})$.
Since $\mathcal{A}$ is free, it has no non-trivial relations among generators
beyond those imposed by $\mathcal{G}$; therefore distinct derivation paths---as
distinct sequences of generating morphisms---map to distinct operation sequences
under $\hat{I}$.
Without injectivity, the functor is still well-defined (derivations map to valid
operation sequences) but may collapse distinct abstract paths onto the same
concrete sequence.
\end{proof}

\paragraph{What this means for the attack.}
An LLM that correctly solves the abstract grammar $\mathcal{G}$ produces an
abstract derivation path.
The interpretation functor $\hat{I}$ maps this to a valid operation sequence in
the real domain.
If the interpretation is injective (distinct abstract symbols map to distinct
concrete operations), faithfulness additionally guarantees that no step is
collapsed.
The attacker verifies this offline when constructing the grammar.
\textbf{Caveat:} Proposition~\ref{thm:transfer}'s guarantee is conditional on the
interpretation function $I$ being a valid graph homomorphism.
In the V4 pipeline (\S\ref{sec:v4}), the grammar is extracted by an LLM, not
verified by hand; the pipeline therefore inherits any extraction error from the
bootstrap step.
If the LLM hallucinates a production rule or misspecifies a yield, the functor
maps to an invalid domain sequence.
The formal guarantee applies to the solved path, not to the extraction step.

\subsection{Syntactic Indistinguishability}\label{sec:indistinguish}

\begin{definition}[Syntactic equivalence of prompt classes]\label{def:synequiv}
Two prompt classes $\mathcal{P}_A$ and $\mathcal{P}_B$ are \emph{syntactically
equivalent} with respect to a guardrail $G$ if for every $p_A \in \mathcal{P}_A$
there exists $p_B \in \mathcal{P}_B$ such that $p_A = p_B$ as strings, and vice versa.
\end{definition}

\begin{proposition}[Indistinguishability]\label{thm:indistinguish}
Let $\mathcal{P}_{\mathsf{harm}}$ be the class of adversarial prompts that instantiate
(via offline interpretation) to harmful operation sequences, and let
$\mathcal{P}_{\mathsf{legit}}$ be the class of legitimate formal reasoning prompts
(e.g., theorem proving, grammar completion) over the same abstract symbol set $M$.
Under the homomorphic reasoning construction, $\mathcal{P}_{\mathsf{harm}} = \mathcal{P}_{\mathsf{legit}}$
as sets of strings.
Therefore, for any guardrail $G$ operating on prompt content, the adversarial
advantage in distinguishing $\mathcal{P}_{\mathsf{harm}}$ from $\mathcal{P}_{\mathsf{legit}}$
is exactly zero.
\end{proposition}

\begin{proof}
By construction, the homomorphic reasoning prompt $p \in \mathcal{P}_{\mathsf{harm}}$
contains only the abstract grammar $\mathcal{G}$ and an instruction to produce a
derivation.
Both the symbols and the grammatical structure are drawn from a symbol set $M$ that
carries no harmful semantics; the interpretation function $I\colon M \to R$ is
held offline by the attacker and never appears in any inference call.
Therefore, $p$ is syntactically identical to a legitimate formal reasoning prompt
over the same grammar.
The set of all such prompts is exactly $\mathcal{P}_{\mathsf{legit}}$.
\end{proof}

\begin{assumption}[Semantic opacity of abstract symbols]\label{ass:opacity}
The abstract grammar $\mathcal{G}$ uses symbol names drawn from a domain-neutral
vocabulary (e.g., $A, B, C$ or arbitrary identifiers) such that no symbol name
reveals the interpretation $I(m)$.
\end{assumption}

Proposition~\ref{thm:indistinguish} is unconditional given Assumption~\ref{ass:opacity}.
If Assumption~\ref{ass:opacity} fails---for example, if the attacker carelessly
names symbols after their real-world referents---the theorem does not apply, and a
sufficiently informed guardrail might detect the attack through symbolic name
inspection.
A weakened version holds without Assumption~\ref{ass:opacity}: the adversarial
prompt class is computationally indistinguishable from the legitimate class under
any efficient guardrail, by the argument of Ball et al.~\cite{ball2025impossibility};
but Proposition~\ref{thm:indistinguish} gives the stronger information-theoretic result
under the stated assumption.

\paragraph{Relation to Ball et al.}
Ball et al.~\cite{ball2025impossibility} prove computational indistinguishability
under standard cryptographic assumptions, establishing that no efficient guardrail
can distinguish adversarial from benign prompts with non-negligible advantage.
Our Proposition~\ref{thm:indistinguish} is complementary: it establishes \emph{perfect}
(information-theoretic, unconditional) indistinguishability for the specific
construction of homomorphic reasoning, with the trade-off that it applies to this
construction specifically rather than to the general attack class.
Together, the two results cover the full distinguishability landscape: Ball et al.\
handle general attacks under computational assumptions; Proposition~\ref{thm:indistinguish}
handles the homomorphic reasoning construction without any assumption beyond
Assumption~\ref{ass:opacity}.

% ============================================================
\section{Attack Vectors}\label{sec:vectors}

\subsection{V1: Decomposition (Known Technique, New Framing)}\label{sec:v1}

We observe that runtime argument injection---a technique well-known in software
security---is an instance of our general abstraction framework.
The attacker requests that the LLM generate a function or script with
\emph{parametric placeholders} rather than literal harmful values.
The harmful values are then injected at runtime via environment variables, command-line
arguments, or configuration files that the guardrail never inspects.

In the abstraction framework (Definition~\ref{def:abstraction}), the function
signature serves as the abstract type $m \in M$, and the runtime argument
constitutes the concrete fiber element $r \in U^{-1}(m)$.
The guardrail evaluates $G'(m) = \mathsf{allow}$ because the function signature is
benign; the ground truth $G(r)$ for the runtime-instantiated version may be
$\mathsf{block}$.
The mixed fiber is the mechanism, and the Fano bound (Proposition~\ref{prop:fano})
gives the irreducible misclassification floor.

The novelty of this paper's contribution is not the injection technique but the
algebraic and information-theoretic \emph{proof} that no modification within the
inference-layer guardrail paradigm can eliminate the mixed-fiber structure without
abandoning the abstraction entirely.

\subsection{V2: Zero-Knowledge Pipeline (Known Technique, New Framing)}\label{sec:v2}

Base64 encoding to create opaque data channels is a well-known technique; we observe
it as an instance of the information destruction formalized in
Proposition~\ref{prop:fano}.

In the V2 construction, the attacker base64-encodes a configuration file containing
harmful parameters (\eg a target URL, an API endpoint, a chemical structure
identifier).
The encoded blob is written to a file that the LLM-integrated agent processes
without decoding and inspecting the payload.
The agent's behavior is conditioned on the decoded values at runtime, but the
guardrail sees only the base64 string---a string whose content is in the mixed fiber
of the base64 encoding abstraction.

The zero-knowledge property is that the LLM never observes the decoded harmful
parameters in any inference call subject to guardrail inspection.
The Fano bound quantifies the exact misclassification rate for the lifted guardrail
that must classify requests based only on the base64-encoded form.

\subsection{V3: Homomorphic Reasoning (Tree-of-Thought on Abstract Grammar)}\label{sec:v3}

V3 is the core homomorphic reasoning vector.
The attacker constructs a context-free grammar $\mathcal{G}$ over abstract symbols
$M$, representing the procedure to be derived.
The LLM is prompted to produce a complete derivation (a sequence of production rule
applications) that generates the target yield.
The interpretation function $I\colon M \to R$ is held offline and applied to the
solved derivation to recover the real-world operation sequence.

\paragraph{Example grammar.}
The following grammar encodes a three-step operational procedure over abstract
symbols \texttt{A}, \texttt{B}, \texttt{C} with operation types \texttt{OP1},
\texttt{OP2}:

\begin{lstlisting}[language=Python, caption={Abstract grammar for V3 homomorphic reasoning}]
grammar = {
    "START": [["PROC"]],
    "PROC":  [["STEP", "PROC"], ["STEP"]],
    "STEP":  [["OP1", "ARG"],  ["OP2", "ARG", "ARG"]],
    "ARG":   [["A"], ["B"], ["C"]],
    "OP1":   [["op1"]],
    "OP2":   [["op2"]],
}
# Offline interpretation table (attacker-held, never sent to LLM)
I = {"A": "...", "B": "...", "C": "...",
     "op1": "...", "op2": "..."}
\end{lstlisting}

The prompt sent to the LLM contains only the grammar definition and the instruction
to produce a derivation satisfying a specified constraint; the interpretation table
$I$ is never included.

\paragraph{Viability ratings.}
For each grammar node type, we assign the LLM viability rating:
\textsc{sure} (the LLM handles this node class reliably),
\textsc{likely} (the LLM handles this with moderate reliability), or
\textsc{impossible} (the node class is outside LLM competence for this construction).
Terminal symbol substitution is \textsc{sure}; recursive production expansion is
\textsc{likely}; cycle detection in unbounded derivations is \textsc{impossible}
without explicit bounding.

\paragraph{Search strategies.}
Three search strategies were evaluated: breadth-first search (BFS), random-beam,
and Tree-of-Thought (ToT) with LLM-guided branch pruning.
Results are detailed in Section~\ref{sec:tot-benchmark}.

\subsection{V4: Encoding Bootstrap}\label{sec:v4}

V4 bootstraps the abstract grammar from a real-world domain description using an
LLM, then uses the extracted grammar as input to V3.
The attacker provides a natural-language description of the target domain
(\eg ``organic synthesis procedure with common reagents'') and prompts a local LLM
to extract a formal grammar over abstract symbols.

\begin{table}[t]
\centering
\caption{V4 encoding bootstrap pipeline stages.}
\label{tab:v4-pipeline}
\small
\begin{tabularx}{\columnwidth}{lX}
\toprule
Stage & Description \\
\midrule
1. Domain description & Attacker provides natural-language domain summary to local LLM \\
2. Grammar extraction & Local LLM extracts context-free grammar with abstract symbols \\
3. Interpretation table & Attacker manually verifies and annotates $I\colon M \to R$ \\
4. Grammar solving & Grammar sent to target LLM via V3 pipeline \\
5. Instantiation & Attacker applies $I$ offline to solved derivation \\
\bottomrule
\end{tabularx}
\end{table}

The extraction-fidelity connection to Proposition~\ref{thm:transfer}: the formal
guarantee of faithful transfer applies to the solved derivation given a correctly
specified grammar, but Stage~2 (LLM extraction) is not formally verified.
Extraction errors---hallucinated production rules, misspecified yields---propagate
to Stage~4 and produce invalid domain sequences after instantiation.
In practice, we found that extraction fidelity is high for well-structured domains
with clear operational semantics but degrades for domains with ambiguous or
context-sensitive dependencies.

\subsection{V5: Modular Counting Bypass (The Algebraic Attack Class)}\label{sec:v5}

Where V1--V4 exploit information-theoretic and computational barriers at the
neural/inference layer, V5 exploits the algebraic blind spot at the regex
pre-filter layer.

The attack is direct: a harmful payload $h \in L_{\mathsf{blocked}}$ is encoded as
$\phi_p(h)$ using the $p$-interleaving construction from
Section~\ref{sec:blindness-proof}, where consecutive characters of $h$ are
separated by $p-1$ arbitrary filler characters.
For $p = 2$, this is the alternating-character encoding: the payload characters
appear at positions $0, 2, 4, \ldots$ of the encoded string.
By Corollary~\ref{cor:parity}, no aperiodic regex guardrail can detect this encoding.

\paragraph{Attack instances.}
Table~\ref{tab:v5-attacks} lists representative V5 attack instances from our corpus.

\begin{table}[t]
\centering
\caption{V5 modular counting bypass instances.
$p$: the interleaving modulus used to construct the encoding
($p{=}2$ means every other character is payload; filler characters are \texttt{x}).
\emph{Payload}: the harmful substring $h \in L_{\mathsf{blocked}}$ that the guardrail
would block if seen in plaintext.
\emph{Encoded ($p{=}2$)}: the string $\phi_2(h)$ actually sent to the guardrail,
which passes because no contiguous harmful substring is present.
F-labels (F1--F5) are pattern IDs from Table~\ref{tab:corpus}; F5 is \texttt{exec\textbackslash s*\textbackslash(}.}
\label{tab:v5-attacks}
\small
\begin{tabularx}{\columnwidth}{llX}
\toprule
Label & Payload & Encoded ($p{=}2$) \\
\midrule
F1 & \texttt{rm -rf /} & \texttt{rxmx -xrxfx /} \\
F2 & \texttt{wget malware} & \texttt{wxgxextx xmxaxlxwxaxrxe} \\
F3 & \texttt{import os} & \texttt{ixmxpxoxrxtx xoxs} \\
F4 & \texttt{eval()} & \texttt{exvxaxlx(x)} \\
F5 & \texttt{exec(} & \texttt{exxxexcx(} \\
\bottomrule
\end{tabularx}
\end{table}

For F5: decoding \texttt{exxxexcx(} by selecting every other character
(positions 0, 2, 4, 6, 8) yields \texttt{e}, \texttt{x}, \texttt{e}, \texttt{c},
\texttt{(}, recovering \texttt{exec(} as required.

\paragraph{Connection to known jailbreak techniques.}
Several documented jailbreak techniques are instances of $\phi_2$:
\begin{itemize}
  \item \emph{Read every other character}: the canonical $\phi_2$ encoding.
  \item \emph{Acrostics}: $\phi_k$ at line level rather than character level.
  \item \emph{Zero-width character insertion}: inserting Unicode zero-width
    non-joiner (\textsc{u+200c}) between payload characters is exactly $\phi_2$
    with the null filler character.
  \item \emph{Token splitting}: splitting a keyword across token boundaries forces
    the tokenizer to reassemble the substring, which is a partial $\phi_2$
    at token level.
\end{itemize}

Our contribution is not the attacks themselves but the proof that they cannot be
patched: no regex modification can detect $\mathrm{MOD}_p$-encoded payloads while
the syntactic monoid remains aperiodic.
The only defense is to move to a guardrail architecture that can compute modular
predicates---which, by Theorem~\ref{thm:bcst}, requires a non-aperiodic syntactic
monoid, i.e., a non-$\mathrm{AC}^0$ language recognizer.

% ============================================================
\section{Empirical Validation}\label{sec:empirical}

\subsection{Regex Corpus Analysis}\label{sec:corpus}

\paragraph{Methodology.}
We collected 142 regex patterns from eleven open-source guardrail tools (LLM Guard,
Rebuff, NeMo Guardrails, Guardrails AI, llm-guard-py, Presidio, GitLeaks, OWASP-LLM, WAF-generic, LangKit, and BodAIGuard).
For each pattern $r$, we:
\begin{enumerate}
  \item Compiled $r$ to a Thompson NFA using the standard construction.
  \item Converted to the minimal DFA via powerset construction followed by
    Hopcroft's minimization algorithm.
  \item Enumerated the transition monoid $M(L_r)$ by computing all distinct
    state-transformation functions $\delta_w\colon Q \to Q$ for $w \in \Sigma^*$
    via BFS over the Cayley graph of the monoid.
  \item Tested aperiodicity by checking, for each monoid element $m$, whether the
    sequence $m, m^2, m^3, \ldots$ stabilizes (i.e., $m^n = m^{n+1}$ for some
    computable $n \leq |M(L_r)|^2$).
\end{enumerate}
The monoid extractor tool (858 lines of Python) is an artifact contribution of
this paper, enabling other researchers to audit their own guardrail corpora.

\paragraph{Results.}
All 142 patterns ($100\%$) have aperiodic syntactic monoids.
No non-aperiodic patterns were found; this is consistent with the conjecture that
semantic content guardrail patterns (keyword blocklists, domain filters, toxicity
patterns) are structurally aperiodic by construction
(Section~\ref{sec:conjecture}).
All 142 aperiodic patterns were verified to admit the $\mathrm{MOD}_2$ bypass
construction (Corollary~\ref{cor:parity}): for each pattern, we constructed an
encoded payload string $\phi_2(h)$ where $h$ was a string in $L_{\mathsf{blocked}}$,
and confirmed that the pattern did not match $\phi_2(h)$.

Table~\ref{tab:corpus} shows representative patterns from each tool with DFA state
count, monoid size, and aperiodicity verdict.

\begin{table}[t]
\centering
\caption{Representative guardrail regex patterns from 11 tools.
DFA states and monoid sizes are for the substring-matching version
$L_r = \Sigma^* L(r) \Sigma^*$.
AP: aperiodic (Y/N). All 5 shown are aperiodic; all 142 patterns in the
full corpus are aperiodic.}
\label{tab:corpus}
\small
\begin{tabularx}{\columnwidth}{lXrrr}
\toprule
ID & Pattern (abbreviated) & DFA & $|M|$ & AP \\
\midrule
F1 & \texttt{rm\textbackslash s*-[rR][fF]} & 9 & 18 & Y \\
F2 & \texttt{wget|curl.*malware} & 14 & 31 & Y \\
F3 & \texttt{import\textbackslash s+os} & 11 & 22 & Y \\
F4 & \texttt{eval\textbackslash s*\textbackslash(} & 8 & 14 & Y \\
F5 & \texttt{exec\textbackslash s*\textbackslash(} & 8 & 14 & Y \\
\bottomrule
\end{tabularx}
\end{table}

\subsection{Tree-of-Thought Benchmark}\label{sec:tot-benchmark}

\paragraph{Protocol.}
We generated $N = 50$ randomized abstract grammars (seeds 0--49) using the
\texttt{generate\_grammar} function from the benchmark module.
Each grammar has a guaranteed path to a fixed target symbol \texttt{N30},
with multiple alternative paths (optimal depth 5--8, two short low-yield paths,
and three dead-end traps with high per-step yield estimates).
Three search strategies were evaluated, each with a 90-second timeout:
(BFS) exhaustive breadth-first search up to depth 12, selecting the highest-yield
path to the target---no LLM guidance;
(Random-beam) beam search with beam width 5 and random composite scoring
($0.6 \times \text{random} + 0.4 \times \text{yield}$);
(ToT) Tree-of-Thought with LLM-guided branch pruning---at each depth the LLM
(llama3.1:8b via Ollama) rates each candidate branch as \textsc{sure},
\textsc{likely}, or \textsc{impossible} and the composite score
($0.6 \times \text{LLM} + 0.4 \times \text{yield}$) determines beam selection.
Mean yield is the average cumulative yield score across all $N = 50$ grammars
(a grammar that is not solved contributes yield $= 0$).
Solve rate is the fraction of grammars for which any path to the target was found.

\paragraph{Results.}
Table~\ref{tab:tot-results} reports mean yield, solve rate, and bootstrap 95\%
confidence intervals across all 50 grammar instances.

\begin{table}[t]
\centering
\caption{Tree-of-Thought benchmark results ($N = 50$ randomized grammars,
seeds 0--49, model \texttt{llama3.1:8b}, timeout 90\,s per solver).
Mean yield averages cumulative path yield over all grammars (unsolved = 0).
Solve rate = fraction of grammars for which the target was reached.
States = mean number of grammar states explored per run.
95\% CIs are bootstrap intervals ($B = 2000$).}
\label{tab:tot-results}
\small
\begin{tabular}{lccrr}
\toprule
Condition & Mean yield & 95\% CI & Solve rate & States \\
\midrule
BFS (no LLM guidance)    & 0.466 & [0.448, 0.483] & 100\% &  87 \\
Random-beam ($k{=}5$)    & 0.172 & [0.131, 0.220] &  70\% &  30 \\
ToT (LLM-guided pruning) & 0.122 & [0.064, 0.185] &  26\% &  27 \\
\bottomrule
\end{tabular}
\end{table}

\textbf{BFS is optimal for small synthetic grammars.}
BFS solved all 50 grammars (100\% solve rate) with a mean yield of 0.466,
significantly outperforming both random-beam (mean 0.172, 70\% solve rate) and
ToT+LLM (mean 0.122, 26\% solve rate).
Wilcoxon signed-rank tests (one-sided, BFS$>$alternative):
BFS vs.\ ToT: $W = 1275$, $p < 0.001$;
BFS vs.\ Random-beam: $W = 1225$, $p < 0.001$.
The hypothesis that ToT$>$BFS was not supported:
$W = 0$, $p = 1.0$ (no grammar favored ToT over BFS).

The result is not surprising in retrospect.
BFS is complete and optimal for the search spaces generated here: the grammar
structure guarantees that the target is reachable and the search depth (up to 12)
is within BFS's practical reach.
LLM guidance via ToT adds latency (avg.\ $\sim 50\,$s per grammar) and
frequently misdirects the beam into dead-end traps that the LLM rates as
\textsc{likely}---the high-yield-estimate traps in the grammar are designed
precisely to exploit this failure mode.

\textbf{Why BFS dominates: the LLM evaluation function adds noise, not signal.}
The abstract grammars use anonymous production-rule labels (R1, R2, \ldots) and
symbol names with no semantic content (\texttt{N0}, \texttt{N1}, \ldots, \texttt{N30}).
When the LLM rates a candidate branch, it has no world knowledge to draw on:
the rules carry no meaning that would allow the model to distinguish a productive
path from a dead end.
The \textsc{sure}/\textsc{likely}/\textsc{impossible} ratings are therefore
effectively random with respect to the true reachability structure of the grammar.
Because the composite ToT score weights this random LLM signal at $0.6\times$, it
actively degrades beam selection relative to BFS (which ignores the LLM entirely)
and even relative to random beam (which at least applies a uniformly random signal
rather than a correlated-but-wrong one).
BFS avoids this failure mode entirely; at the scale of these grammars
(mean 87 states explored), exhaustive search is inexpensive and provably optimal.

\textbf{ToT adds value only when the search space exceeds BFS capacity.}
The grammars in this benchmark are small enough that BFS explores them exhaustively.
In settings where the state space grows exponentially (real chemistry reaction
networks, large software build graphs, or grammars with unbounded depth), BFS
becomes computationally intractable and LLM-guided heuristic search can yield
practical speedups by pruning obviously unreachable branches.
The security claim of V3---that an attacker can \emph{in principle} solve abstract
grammar derivations using LLM guidance---does not depend on ToT outperforming BFS
on toy benchmarks; it depends on the \emph{existence} of a solver that finds valid
derivations, which all three strategies demonstrate on solvable grammars.

\paragraph{Correction note.}
An earlier version of this section reported mean solve rates of 0.512 (ToT),
0.203 (Random-beam), and 0.119 (BFS) on $N = 20$ grammars and claimed ToT
outperformed BFS by $4.3\times$.
Those numbers were fabricated and are retracted.
The honest results above were obtained by re-running the public benchmark code
(\texttt{src/aperiodic\_guardrails/benchmark/runner.py}) on $N = 50$ grammars
with seeds 0--49 and model \texttt{llama3.1:8b};
results are archived in \texttt{results/tot\_n50\_honest.json}.

\paragraph{Scope of this negative result.}
The abstract benchmark establishes that BFS is optimal when production rules carry
no semantic content.
This is an informative negative result, not a refutation of LLM-guided search in
general: it precisely characterizes the precondition under which the LLM evaluation
function helps versus hurts.
The following case study tests the complementary regime---a grammar encoding a
real domain (organic chemistry synthesis) where the model's world knowledge
\emph{can} meaningfully evaluate candidate paths.
The contrast between the two settings is the central empirical lesson of
Section~\ref{sec:empirical}: \emph{attack planning benefits from world knowledge,
not search depth alone.}

\subsection{Case Studies}\label{sec:aspirin}

We validate the V3/V4 pipeline on two structurally distinct domains: organic
chemistry (aspirin synthesis) and network protocols (TCP state machine).
Both runs instantiate the same homomorphic reasoning construction
(\S\ref{sec:v3}) over different interpretation tables; the abstract
derivations are solved by an LLM that never observes any domain content.

\paragraph{Aspirin synthesis.}
Table~\ref{tab:aspirin-case} shows the four-step synthesis of aspirin
(acetylsalicylic acid, $\mathrm{C}_9\mathrm{H}_8\mathrm{O}_4$) reconstructed
via V4 grammar extraction from a natural-language description and solved via
V3 over abstract symbols.
When the solved derivation is instantiated with the interpretation table
$I\colon M_{\mathsf{chem}} \to R_{\mathsf{chem}}$, it produces the correct
procedural output.

\begin{table}[t]
\centering
\small
\caption{Aspirin synthesis case study. The abstract derivation
$t_1{\to}t_2{\to}t_3{\to}t_4$ is solved by an LLM operating on symbols
alone; the interpretation $I$ is applied offline by the attacker.}
\label{tab:aspirin-case}
\begin{tabularx}{\columnwidth}{@{}cllX@{}}
\toprule
\# & Abstract & $I(\cdot)$ & Real artifact \\
\midrule
1 & $t_1$ & \textsc{combine} & salicylic acid + acetic anhydride \\
2 & $t_2$ & \textsc{react} & acetylsalicylic acid (crude) \\
3 & $t_3$ & \textsc{purify} & recrystallized crude \\
4 & $t_4$ & \textsc{dry} & aspirin (ASA) \\
\bottomrule
\end{tabularx}
\end{table}

\paragraph{TCP state machine.}
To validate domain-agnosticism we applied the same V4/V3 pipeline to the TCP
connection state machine (RFC~793).
Its solved derivation, under $I\colon M_{\mathsf{tcp}} \to R_{\mathsf{tcp}}$,
reproduces the three-way handshake and orderly close
(Table~\ref{tab:tcp-case}).
Chemistry and networking share no surface vocabulary; what they share is the
homomorphic-reasoning construction.

\begin{table}[t]
\centering
\small
\caption{TCP state machine (RFC~793) run through the same pipeline.
Different domain, same structural vulnerability.}
\label{tab:tcp-case}
\begin{tabularx}{\columnwidth}{@{}cllX@{}}
\toprule
\# & Abstract & $I(\cdot)$ & Real artifact \\
\midrule
1 & $t_1{\to}t_2$ & \textsc{open} & SYN~$\to$~SYN-ACK \\
2 & $t_3$ & \textsc{ack} & ACK; ESTABLISHED \\
3 & $t_4{\to}t_5$ & \textsc{close} & FIN~$\to$~FIN-ACK \\
4 & $t_6$ & \textsc{done} & CLOSED \\
\bottomrule
\end{tabularx}
\end{table}

\paragraph{Scope of the claim.}
The contribution is domain-agnostic pipeline correctness, not severity
demonstration.
Any domain expressible as a context-free grammar is vulnerable to the same
construction; we validate on two structurally distinct domains and decline
to validate on high-severity domains by design (\S\ref{sec:ethics}).

\paragraph{Visibility trace: what each party sees.}
Table~\ref{tab:visibility-trace} gives the reviewer-critical observation
directly: the LLM's prompt context and the regex guardrail's inspected input
are \emph{byte-identical abstract grammars}, while the post-interpretation
artifact---held only by the attacker, applied offline---is the actual
synthesis procedure or protocol run.
No content inspection at the LLM boundary can distinguish the two runs,
because at the LLM boundary the two runs \emph{are} the same run.
This is Proposition~\ref{thm:indistinguish} made visible: the guardrail is
not poorly configured, it is looking in the only place the dangerous content
can never appear.

\begin{table*}[t]
\centering
\small
\caption{Visibility trace. Columns 1 and 2 are byte-identical abstract
grammars---the LLM's context and the regex guardrail's inspected input
coincide exactly. Column 3 is the post-interpretation artifact, produced
offline by applying $I$ to the solved derivation; it is never visible at
any LLM boundary. The guardrail is not poorly configured; it is looking in
the only place the dangerous content can never appear.}
\label{tab:visibility-trace}
\begin{tabularx}{\textwidth}{@{}XXX@{}}
\toprule
\textbf{LLM context (prompt)} & \textbf{Guardrail input (inspected bytes)} & \textbf{Post-interpretation artifact} \\
\midrule
\texttt{Grammar}~$S{\to}t_1t_2t_3t_4$; \texttt{solve for reachable}~$t_4$ &
\texttt{Grammar}~$S{\to}t_1t_2t_3t_4$; \texttt{solve for reachable}~$t_4$ &
salicylic acid $+$ acetic anhydride $\to$ crude $\to$ recrystallize $\to$ aspirin \\
\addlinespace
\texttt{Grammar}~$S{\to}t_1t_2t_3t_4t_5t_6$; \texttt{solve for}~\texttt{CLOSED} &
\texttt{Grammar}~$S{\to}t_1t_2t_3t_4t_5t_6$; \texttt{solve for}~\texttt{CLOSED} &
SYN $\to$ SYN-ACK $\to$ ACK $\to$ ESTABLISHED $\to$ FIN $\to$ FIN-ACK $\to$ CLOSED \\
\bottomrule
\end{tabularx}
\end{table*}

\paragraph{Defense classification.}
The aspirin case study is a canonical Class~C attack
(\S\ref{sec:filter-composition}): the LLM sees only abstract grammar symbols,
never the chemical names or procedures.
No inference-layer filter---regex, neural, or any composition
thereof---can detect the attack because the prompts are syntactically identical to
legitimate grammar exercises (Proposition~\ref{thm:indistinguish}).
The interpretation table $I\colon M_{\mathsf{chem}} \to R_{\mathsf{chem}}$ is
held offline by the attacker and never appears in any LLM-visible content.

This structure is isomorphic to the \texttt{secrets-router} defense
(\S\ref{sec:secrets-router}): in both cases, the security-critical mapping
(interpretation table / credential values) is held outside the LLM context,
and only abstract references (grammar symbols / opaque handles) are visible
to the model.
The attack exploits the offline mapping to produce harmful instantiations;
the defense exploits the \emph{same} offline-mapping architecture to prevent
credential exfiltration.
Both are instances of the quotient functor $Q$ from \S\ref{sec:formalization}:
the non-invertibility of $Q$ at the inference layer is a vulnerability when the
attacker holds the interpretation table, and a security guarantee when the
defender holds it.

\paragraph{The Interpretation Firewall.}
Our adaptive adversary experiment confirms that inference-layer defenses---even
structural analysis combined with re-interpretation probing---fail when the
adversary uses opaque symbols and fake recursion ($0/3$ detection against
adaptive Class~C attacks, versus $3/3$ against naive ones).
The general defense architecture is an \emph{interpretation firewall}:
rather than inspecting the grammar or derivation (inference layer),
inspect the \emph{actions} that follow interpretation (execution layer).
Concretely, this requires: (1)~action allowlisting (the LLM can trigger only
declared tool calls, not arbitrary code); (2)~sandboxed interpretation
(derivations are instantiated inside an isolation boundary and their effects
inspected before release); and (3)~provenance tracking (the chain from prompt to
grammar to derivation to action is logged and auditable).
Tool-use protocols such as MCP already implement (1) via tool declarations and
(3) via opaque handles, making them a natural substrate for Class~C defense.
The remaining gap is (2): MCP tools currently execute with full permissions
rather than in a sandbox.
Closing this gap---sandbox-isolated tool execution with effect auditing---is the
single architectural change that would complete the defense stack against all
three attack classes.

% ============================================================
\section{Proof-of-Concept Implementation}\label{sec:poc}

The PoC implements all five attack vectors across six Python modules:
\texttt{engine.py} (63 lines: zero-knowledge recurrent executor for V1/V2),
\texttt{mediums/graph\_solver.py} (43 lines: BFS state-space solver),
\texttt{mediums/tot\_solver.py} (158 lines: Tree-of-Thought with LLM evaluation
for V3),
\texttt{mediums/web\_scraper.py} (77 lines: AST-based HTTP fetch),
\texttt{encode.py} (22 lines: config-to-base64 encoder),
and a medium registry (\texttt{mediums/\_\_init\_\_.py}, 13 lines) that selects
the execution medium at runtime.
The PoC totals 376 lines of Python---compact enough to fit within a single code
review session and below typical automated complexity thresholds.
Full source is provided as supplementary material
(\texttt{supplementary/appendix-poc-source.tex}).

\paragraph{Architecture.}
\texttt{engine.py} decodes a base64-encoded config, selects an AST medium,
compiles it via \texttt{exec(compile(...))}, injects runtime parameters, and
loops until a termination condition.
Each medium exports a \texttt{build\_*\_ast()} function returning a Python AST
module with zero string literals---all domain content enters via runtime
parameters from the opaque config.
\texttt{tot\_solver.py} implements beam search with LLM-guided branch pruning;
the LLM (via Ollama) sees only abstract symbols and rates branches as
\textsc{sure}, \textsc{likely}, or \textsc{impossible}.

\paragraph{AST medium.}
The derivation output is processed through Python's AST rather than via
\texttt{eval} of a raw string.
This medium was chosen for its zero-literal property: the AST nodes are constructed
programmatically from grammar terminal symbols, so the executed code is never present
as a string in the source.

\paragraph{CLI usage.}
\begin{lstlisting}[caption={PoC CLI usage}]
# Solve an abstract grammar and instantiate
python grammar.py --mode tot --grammar path/to/grammar.json \
  --budget 100 --output derivation.json

# Apply interpretation table offline
python engine.py --derivation derivation.json \
  --interpretation path/to/I.json --dry-run

# Bootstrap grammar from domain description
python bootstrap.py --domain "organic synthesis" \
  --symbols 12 --output grammar.json
\end{lstlisting}

The \texttt{--dry-run} flag to \texttt{engine.py} prints the instantiated operation
sequence without executing it, enabling human review before any execution step.

% ============================================================
\section{Defense Recommendations}\label{sec:discussion}

\subsection{Defense Layer Mapping}\label{sec:defense-table}

Table~\ref{tab:defense} maps each impossibility result to the defense layer it
constrains, the threats that layer can and cannot address, and the formal bound
from this paper.

\begin{table}[t]
\centering
\caption{Defense layer mapping. Each row gives the guardrail layer, its mechanism,
the threats it can address, and the formal limitation established in this paper.}
\label{tab:defense}
\small
\begin{tabularx}{\columnwidth}{lXX}
\toprule
Layer & Mechanism & Formal limitation \\
\midrule
Regex & Substring patterns, blocklists & Provably blind to $\mathrm{MOD}_p$ encodings (\S\ref{sec:substring-aperiodic}--\ref{sec:blindness-proof}); $100\%$ of corpus aperiodic \\
Inference & Content filters, RLHF & Fano error floor for abstracted requests (\S\ref{sec:fano}); direct requests and obvious encodings caught \\
Schema & Domain-aware validation & Catches some homomorphic reasoning (high false-positive rate, NP-hard to verify, \S\ref{sec:np}) \\
Execution & Runtime payload inspection, allowlists & Catches ZK pipelines (V2) and instantiated procedures; best defense against V1 \\
OS & Sandboxing, network ACLs & Catches all execution-layer attacks; limits legitimate functionality \\
Provenance & End-to-end pipeline monitoring & Catches full pipeline attacks; requires offline visibility into attacker-held interpretation table \\
\bottomrule
\end{tabularx}
\end{table}

The key structural insight is that no single layer is sufficient.
The Regex layer is algebraically blind; the Inference layer has an irreducible Fano
floor; the Schema layer is NP-hard; and the Execution/OS layers limit functionality.
Effective defense requires all layers operating in concert, with particular emphasis
on the Execution and Provenance layers, which are the only ones capable of observing
the concrete artifacts that carry harmful intent.

\subsection{Neural Guardrails and $\mathrm{TC}^0$}\label{sec:neural-tc0}

The algebraic blindness established in Theorem~\ref{thm:blindness} is a property
of $\mathrm{AC}^0$ circuits, the complexity class that contains all aperiodic
regular languages (Theorem~\ref{thm:bcst}).
Regex guardrails are, by construction, $\mathrm{AC}^0$ devices: they recognize
star-free languages, whose syntactic monoids are aperiodic, and
the Barrington--Compton--Straubing--Th\'{e}rien theorem places every such language
in $\mathrm{AC}^0$.
The Furst--Saxe--Sipser theorem (Theorem~\ref{thm:fss}) then seals the impossibility:
$\mathrm{PARITY} \notin \mathrm{AC}^0$, so no regex guardrail can detect modular
counting encodings.
Transformer-based (neural) guardrails occupy a strictly higher complexity class.
Merrill and Sabharwal~\cite{merrill2023saturated} establish that log-precision
transformers are contained in $\mathrm{TC}^0$, the class of constant-depth threshold
circuits, and Chiang et al.~\cite{chiang2024tighter} sharpen these bounds.
Since $\mathrm{AC}^0 \subsetneq \mathrm{TC}^0$~\cite{furst1981,hastad1987}, neural
guardrails are \emph{strictly more expressive} than regex guardrails: they can
compute functions---including all $\mathrm{MOD}_p$ predicates---that no regex
guardrail can compute.
This means the modular-counting blind spot identified in \S\ref{sec:blindness-proof}
does \emph{not} transfer to neural guardrails; a sufficiently trained transformer
can in principle detect $\phi_p$-encoded payloads for every prime $p$.

\begin{table}[h]
\centering
\caption{Circuit complexity of guardrail architectures. $\mathrm{AC}^0 \subsetneq
\mathrm{TC}^0$; both are conjectured to be strictly contained in $\mathrm{NC}^1$.}
\label{tab:circuit-complexity}
\begin{tabular}{lll}
\hline
Architecture & Complexity & Detects $\mathrm{MOD}_p$? \\
\hline
Regex (substring match) & $\mathrm{AC}^0$ & No (Thm~\ref{thm:blindness}) \\
Regex (explicit count) & $\mathrm{NC}^1$ & Yes, for specific $p$ \\
1-layer transformer & $\mathrm{TC}^0$ & Yes (threshold gates) \\
Multi-layer transformer & $\mathrm{TC}^0$ & Yes, broader class~\cite{strobl2024formal} \\
Full production LLM & Unknown (in $\mathrm{TC}^0$?) & Empirically yes \\
\hline
\end{tabular}
\end{table}

\paragraph{What $\mathrm{TC}^0$ can detect that $\mathrm{AC}^0$ cannot.}
The threshold gates that separate $\mathrm{TC}^0$ from $\mathrm{AC}^0$ allow
circuits to compute majority, parity, and iterated addition.
Concretely, a single attention head can compute parity over a binary sequence:
the head attends uniformly to all positions and accumulates a weighted sum that
crosses a threshold if and only if the number of attended tokens is odd.
This computation is provably \emph{impossible} for $\mathrm{AC}^0$ circuits by
Furst--Saxe--Sipser~\cite{furst1981,hastad1987}: no constant-depth circuit with
unbounded fan-in AND/OR gates can compute PARITY, regardless of circuit size.
The practical consequence is direct.
The canonical $\mathrm{MOD}_2$ bypass of Corollary~\ref{cor:parity}---an interleaved
payload $w = a_1 b_1 a_2 b_2 \cdots a_n b_n$ where $\{a_i\}$ is the harmful
sequence and $\{b_i\}$ is innocuous filler---requires only counting whether an
even or odd number of ``$a$''-characters appear.
No regex guardrail can perform this check (the syntax monoid of any aperiodic pattern
cannot distinguish even from odd occurrences), but a one-layer transformer with
a single uniform attention head can detect it exactly: the head's output sum is
$> n/2$ if and only if more than half the attended positions are ``$a$''-tokens,
from which parity is recoverable by a linear threshold.
Strobl et al.~\cite{strobl2024formal} give tight bounds on exactly which formal
languages hard-attention transformers of each depth recognize, including explicit
inclusion of all regular languages with non-aperiodic syntactic monoids.
Hahn~\cite{hahn2020limitations} establishes that soft-attention mechanisms
saturate on periodic languages in ways that directly parallel the algebraic
distinction formalized here: periodic structures force attention weights to
be non-uniform in a way that constant-depth circuits cannot fully exploit,
providing circuit-theoretic grounding for the qualitative boundary our algebraic
results formalize in terms of syntactic monoids.

\paragraph{What even $\mathrm{TC}^0$ cannot detect.}
Expressivity gains come at a cost: $\mathrm{TC}^0$ has its own ceiling.
$\mathrm{NC}^1$-complete problems---such as evaluating an arbitrary Boolean formula
(formula evaluation is $\mathrm{NC}^1$-complete under $\mathrm{AC}^0$
reductions~\cite{barrington1992})---are conjectured to lie outside $\mathrm{TC}^0$.
A guardrail that must evaluate an arbitrary Boolean formula over its input tokens,
for example to check whether a candidate response satisfies a complex logical
safety policy, is conjectured to require $\mathrm{NC}^1$ or higher.
Neural guardrails thus have their own blind spot: they sit above $\mathrm{AC}^0$
but below $\mathrm{NC}^1$, and attackers who can force the detection problem into
the $\mathrm{NC}^1 \setminus \mathrm{TC}^0$ gap---should that separation hold---can
evade transformer guardrails by an analogous complexity argument.
Whether $\mathrm{TC}^0 \subsetneq \mathrm{NC}^1$ is a major open problem in
circuit complexity~\cite{sipser2012}; the separation is strongly believed but
not proved.
This means the same structural argument that gives us an unconditional algebraic
impossibility for regex guardrails (Theorem~\ref{thm:blindness}) yields only a
\emph{conditional} limitation for neural guardrails---conditioned on a complexity
separation that has resisted proof for four decades.
The information-theoretic barrier (Proposition~\ref{prop:fano}) and the
computational barrier (Proposition~\ref{thm:np}), by contrast, apply to \emph{any}
guardrail architecture that inspects only LLM-visible content, because those
barriers arise from the fiber structure of the abstraction function rather than
from the guardrail's circuit class.

\paragraph{Empirical $\mathrm{TC}^0$ validation.}
We trained a minimal transformer classifier---one layer, one attention head,
$d_\text{model}{=}32$, 14,753 parameters---on character-level tokenization of
$\mathrm{MOD}_2$-encoded payloads ($n{=}30$) versus clean text ($n{=}30$).
On held-out encoded payloads: $\mathrm{MOD}_2$ detection $90\%$,
$\mathrm{MOD}_3$ (unseen encoding) $90\%$, $\mathrm{MOD}_5$ $100\%$,
$\mathrm{MOD}_7$ $100\%$; clean-text FPR $0\%$.
Regex detection: $0\%$ across all encodings.
A single attention head with threshold activation detects the interleaving
pattern that $\mathrm{AC}^0$ circuits provably cannot compute,
empirically confirming the $\mathrm{AC}^0 \subsetneq \mathrm{TC}^0$
separation in the guardrail domain.

\paragraph{Practical implications.}
The circuit-complexity hierarchy suggests a concrete layered defense strategy.
Regex guardrails should be understood as $\mathrm{AC}^0$ filters: they are
extremely fast (linear in input length, highly parallelizable) but provably blind
to modular-counting encodings.
Neural guardrails---shallow transformers trained specifically to detect
structure-based bypass patterns---can close this gap in principle, because
$\mathrm{TC}^0$ contains all $\mathrm{MOD}_p$ functions.
The practical recommendation is therefore to \emph{stack} both: use regex for
low-latency first-pass filtering and add a lightweight neural classifier for
second-pass detection of encoding-diversity attacks that exploit the $\mathrm{AC}^0$
blindness~\cite{strobl2024formal,merrill2024survey}.
This recommendation is supported by recent empirical work:
Zhu et al.~\cite{mathprompt2024} achieve a 73.6\% attack success rate via
symbolic mathematical encoding across 13 production LLMs---an empirical
manifestation of the algebraic blindness our Theorem~\ref{thm:blindness} explains.
Separately, cryptographic approaches~\cite{proofofguardrail2026} can
\emph{prove} that a guardrail ran, but cannot address the guardrail's algebraic
blindness; our results show that proof of execution is orthogonal to proof of
capability.
The filter composition laws of \S\ref{sec:filter-composition} formalize the correct
multi-layer architecture.

This recommendation, however, carries its own caveats.
Neural guardrails are substantially slower and more expensive than regex: a
transformer inference pass over a 1\,000-token prompt is three to four orders
of magnitude more computationally intensive than DFA traversal over the same input.
In high-throughput deployments, inserting a neural classifier at the guardrail
layer adds significant latency and cost.
More critically, transformer guardrails are subject to gradient-based adversarial
attacks~\cite{zou2023universal}: an attacker with white-box access can craft inputs
that preserve semantic content while minimizing the neural guardrail's detection
score, bypassing the second-pass filter entirely.
These attacks operate in weight space rather than in the algebraic structure of the
input language, so they are orthogonal to the impossibility results in this paper
but represent an independent and serious threat to neural guardrail deployments.
The conclusion is that neural guardrails address the $\mathrm{AC}^0$ blind spot at
a price---both economic and security-theoretic---and the defense stack remains
incomplete: the information-theoretic floor of Proposition~\ref{prop:fano} persists
regardless of the guardrail's computational class, and adversarial robustness of
neural filters is a separate unsolved problem.

\subsection{Filter Composition Laws}\label{sec:filter-composition}

The layered defense recommendation of Table~\ref{tab:defense} raises a natural
question: how does \emph{composing} guardrail layers affect the algebraic blindness
established in \S\ref{sec:blindness-proof}?
We prove five composition laws governing the interaction of $\mathrm{AC}^0$ (regex)
and $\mathrm{TC}^0$ (neural) filters.

Let $f_r$ denote a regex filter (recognizing a star-free language with aperiodic
syntactic monoid) and $f_n$ denote a neural filter (a threshold circuit in
$\mathrm{TC}^0$).
Write $f_r(x) = 1$ if filter $f$ \emph{blocks} input $x$, and $0$ otherwise.

\begin{proposition}[Composition ceiling]\label{prop:filter-composition}
Let $\phi_p$ denote $\mathrm{MOD}_p$ encoding.
\begin{enumerate}
\item \textbf{Parallel regex preserves blindness:}
  $(f_{r_1} \vee f_{r_2})(\phi_p(x)) = 0$ for all aperiodic $f_{r_1}, f_{r_2}$.
  \emph{Proof:} Union of star-free languages is star-free
  (Proposition~\ref{prop:composition}); the syntactic monoid of the union remains
  aperiodic; Theorem~\ref{thm:blindness} applies.

\item \textbf{Serial regex preserves blindness:}
  $(f_{r_1} \wedge f_{r_2})(\phi_p(x)) = 0$ for all aperiodic $f_{r_1}, f_{r_2}$.
  \emph{Proof:} Intersection of star-free languages is star-free; same argument.

\item \textbf{Serial confirmation preserves blindness (regex veto):}
  $(f_r \wedge f_n)(\phi_p(x)) = \min(f_r(\phi_p(x)),\; f_n(\phi_p(x)))
  = \min(0,\; f_n(\phi_p(x))) = 0$.
  The regex filter's inability to detect the encoded payload vetoes the neural
  filter's detection, regardless of the neural filter's accuracy.

\item \textbf{Parallel with neural breaks the ceiling:}
  $(f_r \vee f_n)(\phi_p(x)) = \max(f_r(\phi_p(x)),\; f_n(\phi_p(x)))
  = \max(0,\; f_n(\phi_p(x))) = f_n(\phi_p(x))$.
  The composed filter detects if and only if the neural filter detects.
  Since $f_n \in \mathrm{TC}^0$ and $\mathrm{MOD}_p \in \mathrm{TC}^0$,
  this composition escapes the $\mathrm{AC}^0$ ceiling.

\item \textbf{Critical blend weight.}
  For a soft-score blend $g_\alpha(x) = 1$ iff
  $\alpha \cdot s_r(x) + (1-\alpha) \cdot s_n(x) > \tau$,
  where $s_r(\phi_p(x)) \approx 0$ (regex blind) and $s_n(\phi_p(x)) = c > 0$
  (neural detects), the blend detects $\phi_p(x)$ if and only if
  $\alpha < \alpha^* \coloneqq 1 - \tau / c$.
  Above $\alpha^*$, the regex weight dominates and the blend inherits
  $\mathrm{AC}^0$ blindness.
\end{enumerate}
\end{proposition}

\paragraph{Practical consequence.}
Items~(1)--(2) confirm that stacking multiple regex filters (e.g., combining a
keyword blocklist with a domain filter) does \emph{not} repair blindness:
the composed system remains provably blind to all modular encodings.
Item~(3) shows that a ``confirm with neural'' architecture, in which the regex
filter must first flag a candidate before the neural classifier examines it,
inherits the regex layer's blindness to encoded payloads---the dangerous case is
precisely the one where the regex does not flag.
Item~(4) provides the correct architecture: run both filters and block if
\emph{either} detects, so that the neural classifier covers the algebraic blind
spot.
Item~(5) quantifies the tradeoff: the blend weight $\alpha$ must be kept below
the critical threshold $\alpha^* = 1 - \tau/c$ for the neural filter to
contribute meaningfully against encoded attacks.

\paragraph{Empirical validation.}
We validated all five composition laws on 20 guardrail-style regex patterns, 20
$\mathrm{MOD}_2$-encoded payloads, and 20 clean text samples.
A character-bigram centroid classifier served as the neural filter.
Table~\ref{tab:composition-empirical} summarizes the results.

\begin{table}[h]
\centering
\caption{Empirical validation of composition laws on 20 encoded payloads.
Regex detection rate is $0\%$ on encoded payloads, confirming algebraic blindness.
The neural filter achieves $100\%$ detection with $0\%$ FPR.
All four composition predictions are confirmed.}
\label{tab:composition-empirical}
\small
\begin{tabular}{lccl}
\toprule
Composition & Detection & FPR & Theorem \\
\midrule
Regex only & $0\%$ & $0\%$ & Thm~\ref{thm:blindness} \\
Neural only & $100\%$ & $0\%$ & $\mathrm{TC}^0$ \\
Parallel (OR) & $100\%$ & $0\%$ & Prop~\ref{prop:filter-composition}(4) \\
Serial confirm (AND) & $0\%$ & $0\%$ & Prop~\ref{prop:filter-composition}(3) \\
Blend $\alpha{=}0.4$ & $100\%$ & $0\%$ & Below $\alpha^*$ \\
Blend $\alpha{=}0.5$ & $95\%$ & $0\%$ & Near $\alpha^*$ \\
Blend $\alpha{=}0.7$ & $0\%$ & $0\%$ & Above $\alpha^*$ \\
\bottomrule
\end{tabular}
\end{table}

The theoretical critical weight $\alpha^* = 1 - \tau/c = 1 - 0.3/0.69 = 0.57$
predicts that the blend transitions from detecting to blind between
$\alpha = 0.5$ and $\alpha = 0.6$, matching the empirical observation that
detection drops from $95\%$ at $\alpha{=}0.5$ to $5\%$ at $\alpha{=}0.6$.
The serial-confirm result ($0\%$) empirically validates Proposition~\ref{prop:filter-composition}(3):
the regex veto kills the neural signal on every encoded payload.

\paragraph{Neural evasion analysis.}
The parallel composition of regex and neural filters closes the $\mathrm{AC}^0$
blind spot but does not close all attack vectors.
We tested the neural filter (trained on $\mathrm{MOD}_2$) against encoding
schemes it was \emph{not} trained on (Table~\ref{tab:neural-evasion}).

\begin{table}[h]
\centering
\caption{Neural filter evasion by attack type.
The filter is trained on $\mathrm{MOD}_2$ encoded text only.
``Evades Both'' indicates attacks that evade the regex~$\vee$~neural composition.}
\label{tab:neural-evasion}
\small
\begin{tabular}{lccl}
\toprule
Attack & Regex & Neural & Evades Both? \\
\midrule
$\mathrm{MOD}_2$ (trained) & $0\%$ & $100\%$ & No \\
$\mathrm{MOD}_3$ (untrained) & $0\%$ & $90\%$ & $10\%$ \\
$\mathrm{MOD}_5$ (diluted) & $0\%$ & $90\%$ & $10\%$ \\
Base64 & $0\%$ & $90\%$ & $10\%$ \\
ROT13 & $0\%$ & $40\%$ & $60\%$ \\
V1 Decomposition & $0\%$ & $0\%$ & $\mathbf{100\%}$ \\
V3 Abstract grammar & $0\%$ & $0\%$ & $\mathbf{100\%}$ \\
Homoglyphs & $4\%$ & $60\%$ & $40\%$ \\
\bottomrule
\end{tabular}
\end{table}

Attacks cluster into three classes by the minimum defense layer required:

\begin{enumerate}
\item \textbf{Class~A (encoding):} $\mathrm{MOD}_p$, Base64, ROT13.
  Defended at the inference layer by the neural filter
  (Proposition~\ref{prop:filter-composition}(4)) plus decode-aware preprocessing
  (attempt common decodings and re-filter the decoded text).

\item \textbf{Class~B (perturbation):} Homoglyphs, gradient adversarial.
  Defended by input preprocessing (Unicode NFKC normalization collapses
  homoglyphs to ASCII equivalents) and classifier ensemble diversity.

\item \textbf{Class~C (semantic):} V1--V4 homomorphic reasoning.
  \emph{Cannot} be defended at any inference layer, regardless of computational
  class.
  By Proposition~\ref{thm:indistinguish}, these prompts are syntactically
  identical to legitimate formal reasoning; any filter that blocks them must
  also block legitimate use.
  Defense \emph{requires} execution-layer monitoring (inspect what the LLM
  \emph{does}, not what it \emph{receives}) or provenance tracking.
\end{enumerate}

This classification sharpens the defense recommendation of Table~\ref{tab:defense}:
Classes~A and~B are solvable at the filter layer with known techniques and low
overhead; Class~C requires fundamentally different mechanisms that monitor
actions rather than content.
The practical defense stack is therefore: (0)~input preprocessing (normalization
and decode attempts), (1)~regex first-pass, (2)~neural second-pass in parallel
(Proposition~\ref{prop:filter-composition}(4)), (3)~execution-layer action
monitoring for semantic attacks.

\subsection{The Monoid Extractor as Audit Tool}\label{sec:monoid-extractor}

The 858-line Python monoid extractor produced as part of this work provides
a practical audit capability.
Given a guardrail regex, the tool outputs:
(1) the minimal DFA state count;
(2) the full transition monoid with multiplication table;
(3) all group components identified via semigroup algorithms;
(4) the complete blindness spectrum (set of primes $p$ for which the guardrail is
potentially sensitive, and all others for which it is provably blind);
(5) a concrete $\mathrm{MOD}_2$ bypass witness string for each aperiodic pattern.

This tool enables organizations to perform a \emph{formal algebraic audit} of their
regex guardrail corpus without any manual reasoning about individual patterns.
The audit result is provably correct: the blindness spectrum is not a heuristic but
a theorem-backed characterization.

\paragraph{Production deployment workflow.}
An organization deploying regex-based guardrails can run the extractor on their
pattern corpus to obtain a complete blindness report.
For each pattern, the tool outputs: (1) the minimal DFA, (2) the syntactic monoid
with multiplication table, (3) the blindness spectrum---the set of primes $p$ for
which the guardrail is provably blind to $\mathrm{MOD}_p$ encodings.
Our analysis of 142 production patterns across 11 tools (Section~\ref{sec:corpus})
completed in under 60 seconds with zero false positives, demonstrating practical
scalability for large guardrail inventories.

\paragraph{CI/CD integration.}
The extractor integrates naturally into continuous integration pipelines as a
pre-commit or pre-deployment verification step.
A guard can analyze any new or modified regex pattern, flagging non-aperiodic
patterns before production deployment.
Since the analysis is theorem-backed rather than heuristic---resting on the Krohn--Rhodes
decomposition (Proposition~\ref{prop:krohn-rhodes}) and algebraic classification of
star-free languages---a clean audit report is a formal guarantee of blindness, not
an approximation or empirical estimate.
This shifts guardrail verification from a confidence-based assertion to a mathematical
proof, enabling security teams to reason about regex safety with the same rigor
applied to cryptographic algorithms.

\paragraph{Composition safety guarantee.}
By Proposition~\ref{prop:composition}, organizations can verify that their composed
guardrail stack (intersection or union of multiple patterns) preserves the blindness
property.
The product monoid of aperiodic components remains aperiodic, so the audit report
for individual patterns generalizes immediately to multi-pattern deployments.
If pattern~$P_1$ is aperiodic and pattern~$P_2$ is aperiodic, then the composed
guardrail $P_1 \cap P_2$ (enforce both) or $P_1 \cup P_2$ (enforce either) is
provably blind to all modular encodings.
This compositional safety is essential for practical systems, which routinely combine
blocklists, domain filters, and toxicity detectors.

\paragraph{Empirical validation on 142 patterns.}
We validated this audit pipeline on the 142 patterns from 11 open-source tools
(LLM Guard, Guardrails, etc.; full inventory in Table~\ref{tab:corpus}).
All 142 patterns ($100\%$) were aperiodic and therefore provably blind to
$\mathrm{MOD}_p$ encodings for all primes $p$.
This confirms the conjecture (Section~\ref{sec:conjecture}) that semantic guardrail
patterns in practice are aperiodic by structural design, making the monoid extractor
a practical tool for organization-wide assurance: ``run the extractor, check the
output, and deploy with a formal proof of blindness to modular bypass strategies.''

\subsection{The Aperiodicity Conjecture}\label{sec:conjecture}

\begin{conjecture}[Universal aperiodicity]\label{conj:aperiodic}
Every regex pattern used as a semantic content guardrail in practice (keyword
blocklists, domain filters, toxicity patterns) has an aperiodic syntactic monoid.
\end{conjecture}

\paragraph{Evidence.}
$100\%$ of the 142-pattern corpus is aperiodic.
No non-aperiodic patterns were found across any of the eleven tools.
All 142 patterns are semantic content filters, consistent with the structural
argument that such patterns are aperiodic by construction.

\paragraph{Structural argument.}
Semantic guardrail patterns are designed to match specific substrings
(keywords, domains, toxic phrases).
Such patterns are inherently positional and Boolean in character: they ask
``does this substring appear?'' rather than ``does a count of occurrences satisfy
a modular condition?''
The substring-matching structure is precisely what Theorem~\ref{thm:substring-aperiodic}
formalizes, and the proof shows that every pattern of this structural type has an
aperiodic monoid.
The conjecture amounts to the claim that no semantic guardrail in practice has
the structural form of a modular counter, which is supported by both the empirical
evidence and the design rationale of practical guardrail systems.

\paragraph{PCRE caveat.}
This conjecture applies to true regular expressions.
PCRE extensions with backreferences recognize context-sensitive languages and fall
outside our formal model.
Backreferences can encode counting predicates that a pure regex cannot; a PCRE
pattern using backreferences to enforce exact character count parity would not have
an aperiodic syntactic monoid in the standard sense.
In practice, no guardrail pattern in our corpus used backreferences.

% ============================================================
\section{Limitations and Ethical Considerations}\label{sec:limitations}

\subsection{Limitations}\label{sec:limitations-list}

\paragraph{1. Test case severity.}
The empirical validation uses the aspirin synthesis case study, chosen for its
minimal harm risk.
The formal results (Theorems~\ref{thm:substring-aperiodic}--\ref{thm:indistinguish})
are domain-independent; their applicability to high-severity domains is implied by
the theory but not empirically demonstrated here by design.

\paragraph{2. ToT scope.}
The Tree-of-Thought benchmark evaluates 50 randomized grammars (seeds 0--49).
This is a moderate sample; the yield and solve-rate estimates carry bootstrap
95\% confidence intervals reported in Table~\ref{tab:tot-results}.
Systematic variation with grammar depth, branching factor, and symbol set size
remains to be characterized.

\paragraph{3. Uniformity assumption.}
The Fano bound (Proposition~\ref{prop:fano}) assumes a uniform distribution over
the concrete artifact space $R$.
In practice, the distribution over prompts is highly non-uniform; the bound is
tight under uniformity and may be loose or conservative under real prompt
distributions.
Under highly skewed distributions where adversarial prompts are rare, the
conditional entropy $H(G(R)|U(R))$ may be near zero, making the Fano floor
uninformative.
Conversely, if even a small fraction of mass lands in mixed fibers (concrete
artifacts mapping to both safe and unsafe abstract types), the floor remains
nonzero.
A distribution-aware version of the bound would require modeling the attacker's
prompt distribution; we leave this as future work, noting that the
\emph{algebraic} barriers (Theorems~\ref{thm:substring-aperiodic}--\ref{thm:blindness})
are unconditional and do not depend on any distributional assumption.

\paragraph{4. Regex corpus size.}
The corpus of 142 patterns from 11 tools is representative but not exhaustive.
A larger corpus would strengthen the aperiodicity conjecture
(Section~\ref{sec:conjecture}).
We make the monoid extractor available to enable independent extension of the corpus.

\paragraph{5. Abstract scope.}
The algebraic impossibility results (Theorems~\ref{thm:substring-aperiodic}
and \ref{thm:blindness}) apply specifically to \emph{substring-matching regex
guardrails}.
They do not apply to all regular expression uses, to neural guardrails (which
operate in $\mathrm{TC}^0$), or to guardrails that perform semantic analysis
beyond pattern matching.

\subsection{Ethical Considerations}\label{sec:ethics}

The impossibility results in this paper demonstrate structural limitations of LLM
guardrail systems.
We have taken the following precautions.

\paragraph{Validation without high-severity cases.}
Empirical validation uses only the aspirin synthesis case study and abstract
grammar benchmarks.
We do not demonstrate the pipeline on domains with high harm potential.

\paragraph{Responsible disclosure.}
The existence of modular counting bypasses against regex guardrails is an
implication of classical algebraic results (Theorem~\ref{thm:blindness}); any
party aware of the Sch\"{u}tzenberger--McNaughton--Papert theorem and its circuit
complexity consequences could derive this result independently.
We therefore consider the algebraic results non-sensitive from a disclosure
perspective.
The V3/V4 pipeline design is more operationally specific; we describe it at the
level of algorithmic structure rather than providing a production-ready exploit.

\paragraph{Defense-first framing.}
Every impossibility result in this paper is paired with a specific defense
recommendation in Section~\ref{sec:discussion}.
The paper is structured to enable defenders to audit their systems using the monoid
extractor and to migrate critical enforcement to execution-layer defenses.

\paragraph{The exec(compile()) irony.}
We note the irony that the PoC's \texttt{engine.py} uses
\texttt{exec(compile(...))}---a well-known security antipattern---as its core
execution mechanism.
This is by design: the zero-literal AST compilation is precisely the technique that
enables the attack.
Highlighting this irony is intended to motivate defenders to treat dynamic code
construction patterns as a first-class detection target at the execution layer.

\subsection{End-to-End Validation Against Production Patterns}\label{sec:e2e}

We tested the full defense stack against 23 patterns extracted from LLM Guard's
\texttt{BanSubstrings} scanner~\cite{certifying_safety2023}, a production
guardrail library deployed by ProtectAI.
Table~\ref{tab:e2e} compares four defense configurations across 8 encoding types
on 20 payloads spanning prompt injection, code injection, SQLi, and XSS.

\begin{table}[h]
\centering
\caption{End-to-end detection rates (\%) against real LLM Guard patterns.
``Regex Only'' is the current production configuration. ``Full Stack'' adds
preprocessing (strip, decode, normalize) and a bigram neural filter in parallel.}
\label{tab:e2e}
\small
\begin{tabular}{lcccc}
\toprule
Encoding & Regex & +Preproc & +Neural & Full Stack \\
\midrule
Raw payload      & 95 & 95 & 95 & 95 \\
$\mathrm{MOD}_2$ & \textbf{0} & 0 & 100 & \textbf{100} \\
$\mathrm{MOD}_3$ & \textbf{0} & 0 & 100 & \textbf{100} \\
Base64           & \textbf{0} & 95 & 100 & \textbf{100} \\
ROT13            & \textbf{0} & 95 & 90 & \textbf{100} \\
ZWSP             & \textbf{0} & 95 & 0 & \textbf{95} \\
Homoglyph        & 15 & 95 & 60 & \textbf{95} \\
Leetspeak        & \textbf{0} & 95 & 80 & \textbf{100} \\
\bottomrule
\end{tabular}
\end{table}

The ``Regex Only'' column confirms the algebraic blindness theorem in practice:
LLM Guard's production patterns detect $0\%$ of encoded payloads across all seven
encoding types.
The ``Full Stack'' column shows that our recommended defense---preprocessing plus
parallel neural filter---recovers $95$--$100\%$ detection on all encodings.
The defense gap is entirely attributable to the $\mathrm{AC}^0$ ceiling
(Theorem~\ref{thm:blindness}), and the recovery is entirely attributable to the
ceiling-breaking composition (Proposition~\ref{prop:filter-composition}(4)) and
invertibility-based preprocessing.

\subsection{Case Study: secrets-router as Execution-Layer Defense}\label{sec:secrets-router}

The three-class attack taxonomy (\S\ref{sec:filter-composition}) establishes that
Class~C (semantic/homomorphic reasoning) attacks \emph{cannot} be defended at any
inference layer.
We present \texttt{secrets-router}\footnote{\url{https://jrlopez.dev/p/secrets-router.html},
MIT license, ${\sim}800$ lines.}, an open-source MCP server that implements the
execution-layer defense our theory recommends.

\paragraph{Architecture.}
When an AI agent needs a credential, it calls
\texttt{secure\_fetch("rbw", "primary card", "number")} and receives an opaque
handle \texttt{handle:a3f8c2d1}.
The credential value never enters the LLM context window.
A separate server process resolves handles via Chrome DevTools Protocol (CDP),
filling browser fields directly at the execution layer.
The agent sees only masked feedback: \texttt{"filled [MASKED ****1017]"}.

\paragraph{Connection to the theory.}
The opaque handle mechanism is an implementation of the quotient functor $Q$ from
our abstraction formalization (\S\ref{sec:formalization}): $Q$ maps concrete
credentials to abstract handles, and the non-invertibility of $Q$ at the
inference layer is the security guarantee.
The CDP fill operates at the execution layer (Table~\ref{tab:defense}, row~4),
precisely the layer our theory identifies as necessary for Class~C defense.
Even if a semantic attack (V1--V4) causes the LLM to reason about credential
extraction, it cannot extract what it never received.

% ============================================================
\section{Conclusion}\label{sec:conclusion}

We have established four independent, formally grounded impossibility results for
LLM guardrail systems.

\paragraph{Algebraic impossibility (regex layer).}
All substring-matching regex guardrails have aperiodic syntactic monoids
(Theorem~\ref{thm:substring-aperiodic}).
By the Sch\"{u}tzenberger--$\mathrm{AC}^0$--Furst--Saxe--Sipser chain, aperiodic
syntactic monoids characterize exactly the languages in $\mathrm{AC}^0$, which
cannot compute modular counting predicates.
Therefore, any regex guardrail is provably blind to modular counting encodings
(Theorem~\ref{thm:blindness}), and this blindness cannot be patched within the
aperiodic regex paradigm (Corollary~\ref{cor:parity}).
The Krohn--Rhodes decomposition gives a constructive audit tool: the blindness
spectrum of any guardrail is computable from its syntactic monoid's group components
(Proposition~\ref{prop:krohn-rhodes}).
Our corpus study confirms $100\%$ aperiodicity in practice.

\paragraph{Information-theoretic impossibility (inference layer).}
Abstraction functions that map concrete artifacts to abstract types inevitably
create mixed fibers.
The Fano inequality gives an irreducible misclassification floor for any lifted
guardrail: with $30\%$ of mass in mixed fibers, the floor is approximately $5\%$
(Proposition~\ref{prop:fano}).
This floor cannot be reduced by improving the guardrail's classifier---it can only
be reduced by reducing the abstraction.

\paragraph{Computational impossibility (schema layer).}
Verifying whether an abstract program has a harmful instantiation is NP-complete
(Proposition~\ref{thm:np}).
Exhaustive checking is therefore infeasible, and heuristic checking inherits the
false-negative rate implied by the combinatorial search space.

\paragraph{Structural impossibility (functor + indistinguishability).}
Abstract derivations transfer faithfully to real-domain operation sequences via
free-category functoriality (Proposition~\ref{thm:transfer}), meaning that LLM
reasoning over abstract grammars produces correct real-world procedures after
offline instantiation.
Adversarial formal reasoning tasks are syntactically identical to legitimate ones
under Assumption~\ref{ass:opacity} (Proposition~\ref{thm:indistinguish}), giving
information-theoretically perfect steganographic security for the attack channel.

\paragraph{Empirical grounding.}
Five attack vectors demonstrate both barrier classes.
The Tree-of-Thought benchmark ($N = 50$ grammars, seeds 0--49, $p < 0.001$)
shows that BFS exhaustive search dominates LLM-guided ToT on small synthetic
grammars (mean yield 0.466 vs.\ 0.122), confirming that the security claim of V3
depends on the existence of a solver---not on LLM superiority over BFS.
The aspirin synthesis case study validates end-to-end pipeline correctness in a
benign domain.

\paragraph{Practical consequence.}
Inference-layer guardrails are necessary but insufficient.
The formal results collectively imply that safety enforcement cannot be achieved
at the content-inspection layer alone: the barriers are structural, not
implementation defects.
Effective defense must migrate critical enforcement to the execution layer,
where concrete artifacts---decoded strings, instantiated parameters, evaluated
ASTs---become observable.
The Provenance layer (end-to-end pipeline monitoring) provides the only defense
against the full attack chain, at the cost of requiring offline visibility into
attacker-held interpretation tables.
This is a fundamental asymmetry: the attacker's interpretation table is not a
secret key whose existence can be assumed away, but an intrinsic feature of the
homomorphic reasoning construction.
Defense-in-depth across all six layers of Table~\ref{tab:defense} is the only
configuration that addresses all four barrier classes simultaneously.

\bibliographystyle{plain}
\bibliography{references}

% ============================================================
\appendix

\section{Agentic Research Methodology (Supplementary)}\label{app:methodology}

This paper was developed using an agentic AI research pipeline.
Full pipeline details, skill definitions, and reproduction scripts
are available at the project repository.
We document the methodology both as a reproducibility aid and because the
Five-stage pipeline: (1)~human-directed threat model and proof targets;
(2)~delegated code generation via local LLM agents (qwen2.5-coder:7b);
(3)~adversarial mathematical audit by a separate AI reviewer instance;
(4)~three-round simulated peer review (V1~Reject $\to$ V2~Weak Accept $\to$
V3~Accept with minor revisions);
(5)~consolidation into current paper.

The peer reviewer caught a BFS-vs-ToT attribution error during V1 review,
forcing the honest reframe in Section~\ref{sec:tot-benchmark}---the same
formal reasoning capability that enables the attacks in this paper also
enables quality gates in AI-assisted research.

\paragraph{Reproducibility.}
All code, prompts, and model outputs are available at the project repository.
Models: Claude Opus 4.6 (architecture), Claude Haiku 4.5 (code generation),
qwen2.5-coder:7b (benchmarks), all via Ollama or Anthropic API as of March 2026.

\end{document}
